\documentclass[11pt,a4paper,oneside]{report}
\usepackage[T1]{fontenc}
\usepackage[utf8]{inputenc}
\usepackage[italian]{babel}
\usepackage{lmodern}
\usepackage{amsmath,amssymb,amsthm,mathtools}
\usepackage[a4paper,margin=27mm,headheight=14pt]{geometry}
\usepackage{microtype}
\usepackage{enumitem}
\usepackage{booktabs}
\usepackage{needspace}
\usepackage{tikz}
\usetikzlibrary{arrows.meta,calc,patterns}
\usepackage{mdframed}
\usepackage{fancyhdr}
\usepackage{hyperref}
\usepackage{bookmark}
\usepackage{footnotehyper}

\definecolor{coloreschema}{RGB}{30,77,112}
\hypersetup{
  unicode=true,
  hidelinks,
  pdftitle={Appunti di algebra 1},
  pdfauthor={Jude De Sa'Pinto},
  pdfsubject={Algebra 1 -- Prof.ssa E. Detomi},
  bookmarksnumbered=true,
  bookmarksdepth=2
}
% Indice stampato sintetico; sottosezioni disponibili nei segnalibri PDF.
\setcounter{tocdepth}{1}
\setcounter{secnumdepth}{2}
\setlength{\parindent}{0pt}
% Un poco più di aria tra il corpo del testo e il separatore delle note.
\addtolength{\skip\footins}{3pt}
\setlist[enumerate]{leftmargin=*,labelindent=0pt,itemsep=.35em,topsep=.4em}
\setlist[itemize]{leftmargin=*,labelindent=0pt,itemsep=.25em,topsep=.4em}
\setlist[description]{leftmargin=0pt,labelindent=0pt}
\raggedbottom
\interfootnotelinepenalty=10000
\emergencystretch=1.5em
\AddToHook{cmd/section/before}{\Needspace{9\baselineskip}}
\AddToHook{cmd/subsection/before}{\Needspace{7\baselineskip}}

\pagestyle{fancy}
\fancyhf{}
\fancyhead[L]{\small\nouppercase{\leftmark}}
\fancyfoot[C]{\thepage}
\renewcommand{\headrulewidth}{.3pt}
\renewcommand{\chaptermark}[1]{\markboth{#1}{}}

% === STILE MATEMATICO: parametri modificabili ===
% Colori delle barre laterali; titoli e testo restano neri.
\definecolor{coloreRisultati}{RGB}{83,76,153}     % Teoremi, lemmi, proposizioni, corollari, assiomi.
\definecolor{coloreDefinizioni}{RGB}{35,112,77}   % Verde.
\definecolor{coloreEsempi}{RGB}{225,178,24}       % Giallo oro della barra.
\definecolor{coloreEsercizi}{RGB}{202,24,113}     % Rosa acceso.
\definecolor{coloreOsservazioni}{RGB}{90,96,105}  % Grigio neutro.
\definecolor{coloreFasi}{RGB}{130,134,143}        % Contorno delle formule guida.
\definecolor{coloreRichiamiNote}{RGB}{0,85,255}   % Blu elettrico dei numeri delle note.
\newlength{\spessoreBarraAmbiente}
\setlength{\spessoreBarraAmbiente}{1.6pt}
\newlength{\distanzaBarraAmbiente}
\setlength{\distanzaBarraAmbiente}{8pt}
\newlength{\sogliaBloccoBreve}
\setlength{\sogliaBloccoBreve}{8\baselineskip} % I blocchi brevi restano interi.
\newlength{\rientroDimostrazione}
\setlength{\rientroDimostrazione}{4mm}
\newlength{\spazioParagrafiProva}
\setlength{\spazioParagrafiProva}{.42\baselineskip}
\newlength{\spazioFormuleProva}
\setlength{\spazioFormuleProva}{8pt}
\newlength{\spessoreBoxFase}
\setlength{\spessoreBoxFase}{.35pt}
\newlength{\margineBoxFase}
\setlength{\margineBoxFase}{2.2pt}
% Ambienti matematici: una sola barra a sinistra, senza fondo o cornice.
% Le note conservano numerazione, collegamenti e posizione a piè di pagina.
\makeatletter
\renewcommand{\mdf@footnoteinput}{}
% Evita frammenti finali di poche righe; i blocchi lunghi restano spezzabili.
\pretocmd{\detected@mdf@put@frame}{%
  \ifdim\dimexpr\ht\mdf@splitbox@one+\dp\mdf@splitbox@one\relax<\sogliaBloccoBreve
    \booltrue{mdf@nobreak}%
  \fi
}{}{\PackageError{dispense}{Impossibile impostare i blocchi brevi}{Verificare la versione di mdframed.}}
% Colora solo il richiamo, sia nel corpo sia a piè di pagina.
\renewcommand{\@makefnmark}{%
  \hbox{\@textsuperscript{\normalfont\color{coloreRichiamiNote}\@thefnmark}}%
}
\makeatother
\makesavenoteenv{mdframed}
\mdfdefinestyle{barraDispensa}{
  topline=false, rightline=false, bottomline=false, leftline=true,
  linewidth=\spessoreBarraAmbiente,
  backgroundcolor=white, fontcolor=black,
  % La barra è allineata al testo normale; il contenuto rientra di spessore + distanza.
  leftmargin=0pt,
  rightmargin=0pt,
  innerleftmargin=\distanzaBarraAmbiente, innerrightmargin=0pt,
  innertopmargin=1.5pt, innerbottommargin=1.5pt,
  skipabove=6pt, skipbelow=4pt,
  splittopskip=5pt, splitbottomskip=3pt,
  needspace=4\baselineskip,
  nobreak=false
}

\newtheoremstyle{risultatoDispensa}
  {0pt}{0pt}{\itshape}{}{\bfseries}{.}{.55em}
  {\thmname{#1}\thmnumber{ #2}\thmnote{\normalfont\;(#3)}}
\newtheoremstyle{definizioneDispensa}
  {0pt}{0pt}{\normalfont}{}{\bfseries}{.}{.55em}
  {\thmname{#1}\thmnumber{ #2}\thmnote{\normalfont\;(#3)}}
\newtheoremstyle{esempioDispensa}
  {0pt}{0pt}{\normalfont}{}{\bfseries}{:}{.55em}
  {\thmname{#1}\thmnumber{ #2}\thmnote{\normalfont\;(#3)}}
\newtheoremstyle{esercizioDispensa}
  {0pt}{0pt}{\normalfont}{}{\bfseries}{.}{.55em}
  {\thmname{#1}\thmnumber{ #2}\thmnote{\normalfont\;(#3)}}
\newtheoremstyle{osservazioneDispensa}
  {0pt}{0pt}{\normalfont}{}{\bfseries}{.}{.55em}
  {\thmname{#1}\thmnumber{ #2}\thmnote{\normalfont\;(#3)}}

\theoremstyle{risultatoDispensa}
\newtheorem{teorema}{Teorema}[chapter]
\newtheorem{lemma}[teorema]{Lemma}
\newtheorem{proposizione}[teorema]{Proposizione}
\newtheorem{corollario}[teorema]{Corollario}
\newtheorem{assioma}[teorema]{Assioma}
\theoremstyle{definizioneDispensa}
\newtheorem{definizione}[teorema]{Definizione}
\theoremstyle{esempioDispensa}
\newtheorem*{esempio}{Esempio}
\theoremstyle{esercizioDispensa}
\newtheorem*{esercizio}{Esercizio}
\theoremstyle{osservazioneDispensa}
\newtheorem*{osservazione}{Osservazione}

\surroundwithmdframed[style=barraDispensa,linecolor=coloreRisultati]{teorema}
\surroundwithmdframed[style=barraDispensa,linecolor=coloreRisultati]{lemma}
\surroundwithmdframed[style=barraDispensa,linecolor=coloreRisultati]{proposizione}
\surroundwithmdframed[style=barraDispensa,linecolor=coloreRisultati]{corollario}
\surroundwithmdframed[style=barraDispensa,linecolor=coloreRisultati]{assioma}
\surroundwithmdframed[style=barraDispensa,linecolor=coloreDefinizioni]{definizione}
\surroundwithmdframed[style=barraDispensa,linecolor=coloreEsempi]{esempio}
\surroundwithmdframed[style=barraDispensa,linecolor=coloreEsercizi]{esercizio}
\surroundwithmdframed[style=barraDispensa,linecolor=coloreOsservazioni]{osservazione}

% Formule guida in linea: \passomath{A\subseteq B}, \passomath{\Rightarrow}.
% La scatola non altera le formule né introduce numerazione o segnalibri.
\newcommand{\passomath}[1]{%
  \begingroup
  \setlength{\fboxsep}{\margineBoxFase}%
  \setlength{\fboxrule}{\spessoreBoxFase}%
  \fcolorbox{coloreFasi}{white}{\normalfont\ensuremath{#1}}%
  \endgroup\enspace\ignorespaces
}

% Una lista senza etichetta rientra l'intera prova, comprese formule ed elenchi.
% Mantiene le interruzioni di pagina e le note nel normale flusso del documento.
\makeatletter
\renewenvironment{proof}[1][\proofname]{%
  \par\pushQED{\qed}\normalfont
  \list{}{%
    \setlength{\leftmargin}{\rientroDimostrazione}%
    \setlength{\rightmargin}{0pt}%
    \setlength{\labelwidth}{0pt}%
    \setlength{\labelsep}{0pt}%
    \setlength{\itemindent}{0pt}%
    \setlength{\listparindent}{0pt}%
    \setlength{\topsep}{6pt}%
    \setlength{\partopsep}{0pt}%
    \setlength{\parsep}{\spazioParagrafiProva}%
    \setlength{\itemsep}{0pt}%
  }%
  \item\relax
  \setlength{\parskip}{\spazioParagrafiProva}%
  \setlength{\abovedisplayskip}{\spazioFormuleProva}%
  \setlength{\belowdisplayskip}{\spazioFormuleProva}%
  \setlength{\belowdisplayshortskip}{5pt}%
  \textit{#1\@addpunct{.}}\enspace\ignorespaces
}{%
  \popQED\endlist\@endpefalse
}
\makeatother

% Gli svolgimenti usano lo stesso rientro delle dimostrazioni.
\newenvironment{svolgimento}{\begin{proof}[Svolgimento]}{\end{proof}}

% Non lasciare un enunciato isolato in fondo pagina.
\AddToHook{env/proposizione/before}{\Needspace{6\baselineskip}}
\AddToHook{env/teorema/before}{\Needspace{6\baselineskip}}
\AddToHook{env/lemma/before}{\Needspace{6\baselineskip}}
\AddToHook{env/corollario/before}{\Needspace{6\baselineskip}}
\AddToHook{env/assioma/before}{\Needspace{6\baselineskip}}
% === FINE STILE MATEMATICO ===

\newcommand{\N}{\mathbb{N}}
\newcommand{\Z}{\mathbb{Z}}
\newcommand{\Q}{\mathbb{Q}}
\newcommand{\R}{\mathbb{R}}
\newcommand{\C}{\mathbb{C}}
\newcommand{\K}{\mathbb{K}}
\newcommand{\fonti}[1]{\par\nobreak\noindent{\small\itshape #1.}\par\medskip}
\newenvironment{schema}
  {\par\medskip\noindent\begin{minipage}{\linewidth}\centering\small}
  {\end{minipage}\par\medskip}
\tikzset{
  asse/.style={-{Stealth[length=2mm]},draw=black!55,line width=.5pt},
  vettore/.style={-{Stealth[length=2.5mm]},draw=coloreschema,line width=1pt},
  spazio/.style={draw=coloreschema,line width=.9pt,fill=coloreschema!5}
}


\title{Appunti di algebra 1\\[.8em]\large Prof.ssa E. Detomi}
\author{Jude De Sa'Pinto}
\date{}

\begin{document}
\pagenumbering{roman}
\hypersetup{pageanchor=false}
\maketitle
\hypersetup{pageanchor=true}
\tableofcontents
\clearpage

\chapter*{Convenzioni e notazioni}
\addcontentsline{toc}{chapter}{Convenzioni e notazioni}
Scriviamo $\N=\{0,1,2,\ldots\}$, $\N^*=\N\setminus\{0\}$ e $\Z^*=\Z\setminus\{0\}$. Il simbolo $\subseteq$ ammette l'uguaglianza, mentre $\subsetneq$ indica l'inclusione stretta. Il complementare di un insieme è sempre relativo a un ambiente specificato. Indichiamo con $\mathcal P(X)$ l'insieme delle parti di $X$.

La composizione $g\circ f$ applica prima $f$, poi $g$. Per $f:X\to Y$ usiamo $f_*$ per l'immagine diretta e $f^*$ per la controimmagine; quest'ultima esiste anche quando $f$ non è invertibile. La notazione $f^{-1}$, riferita a una funzione, indica l'inversa di una biezione.

Assumiamo le proprietà usuali dei numeri naturali e interi, il principio d'induzione e il buon ordinamento di $\N$. Per i reali assumiamo la struttura di campo ordinato completo e la densità di $\Q$. Le scelte simultanee di rappresentanti, quando necessarie per insiemi arbitrari, sono intese nel consueto quadro con assioma della scelta; i punti in cui intervengono sono precisati nel testo.

\clearpage
\pagenumbering{arabic}
% !TeX root = ../main.tex
\chapter{Insiemi e operazioni}
\fonti{Lezioni 1, 2 e 3}

\section{Operazioni e appartenenza}
% A–L L1 pp. 1–2; M–Z L1. La logica resta proposizionale, senza funzioni caratteristiche.
Un insieme è determinato dai suoi elementi: $A=B$ significa che, per ogni $x$, $x\in A$ se e solo se $x\in B$. Equivalentemente si provano le due inclusioni $A\subseteq B$ e $B\subseteq A$. Si può descrivere un insieme per elenco, come $\{2,3,4\}$, oppure mediante una proprietà: $\Z=\{\ldots,-2,-1,0,1,2,\ldots\}$ e $\Q=\{a/b\mid a,b\in\Z,\ b\ne0\}$.

\begin{definizione}
Per due insiemi $A,B$ si definiscono
\[
\begin{aligned}
A\cup B&=\{x\mid x\in A\ \text{oppure}\ x\in B\},\\
A\cap B&=\{x\mid x\in A\ \text{e}\ x\in B\},\\
A\setminus B&=\{x\mid x\in A\ \text{e}\ x\notin B\},\\
A\triangle B&=(A\setminus B)\cup(B\setminus A).
\end{aligned}
\]
Sono, rispettivamente, \textbf{unione}, \textbf{intersezione}, \textbf{differenza} e \textbf{differenza simmetrica}. L'«oppure» dell'unione ammette che entrambe le appartenenze siano vere; nella differenza simmetrica ne deve essere vera esattamente una. Gli insiemi sono \textbf{disgiunti} se $A\cap B=\varnothing$.
\end{definizione}

\begin{schema}
\begin{tikzpicture}[scale=.67]
  \foreach \t/\lab in {0/Unione,4.5/Intersezione,9/Differenza,13.5/Diff. simmetrica}{
    \begin{scope}[shift={(\t,0)}]
      \ifdim\t pt=0pt\relax\fill[coloreschema!18] (0,0) circle (1);\fill[coloreschema!18] (1.2,0) circle (1);\fi
      \ifdim\t pt=4.5pt\relax\begin{scope}\clip (0,0) circle (1);\fill[coloreschema!18] (1.2,0) circle (1);\end{scope}\fi
      \ifdim\t pt=9pt\relax\begin{scope}\clip (0,0) circle (1);\fill[coloreschema!18] (-1,-1) rectangle (2.3,1);\fill[white] (1.2,0) circle (1);\end{scope}\fi
      \ifdim\t pt=13.5pt\relax\fill[coloreschema!18,even odd rule] (0,0) circle (1) (1.2,0) circle (1);\fi
      \draw (0,0) circle (1);\draw (1.2,0) circle (1);
      \node at (-.35,.5) {$A$};\node at (1.55,.5) {$B$};\node at (.6,-1.5) {\lab};
    \end{scope}
  }
\end{tikzpicture}
\end{schema}

\begin{proposizione}
Unione e intersezione sono commutative e associative:
\[
\begin{gathered}
A\cup B=B\cup A,\qquad A\cap B=B\cap A,\\
(A\cup B)\cup C=A\cup(B\cup C),\qquad
(A\cap B)\cap C=A\cap(B\cap C).
\end{gathered}
\]
Inoltre $A\triangle B=(A\cup B)\setminus(A\cap B)=B\triangle A$.
\end{proposizione}
\begin{proof}
Scambiare le due proposizioni «$x\in A$» e «$x\in B$» non cambia né la loro disgiunzione né la loro congiunzione: si ottiene la commutatività. Per l'associatività, entrambi i membri della prima identità richiedono che $x$ appartenga ad almeno uno tra $A,B,C$; entrambi i membri della seconda richiedono che appartenga a tutti e tre.

Infine, appartenere a uno solo fra $A$ e $B$ equivale ad appartenere alla loro unione, ma non alla loro intersezione. Anche questa condizione è invariata scambiando $A$ e $B$.
\end{proof}

\subsection{Unioni e intersezioni di famiglie}
\begin{definizione}
Sia $(A_i)_{i\in I}$ una famiglia di insiemi, con insieme degli indici $I$. Si pone
\[
\bigcup_{i\in I}A_i=\{x\mid \exists i\in I:\ x\in A_i\},\qquad
\bigcap_{i\in I}A_i=\{x\mid \forall i\in I:\ x\in A_i\}.
\]
Per l'intersezione si suppone $I\ne\varnothing$, oppure si fissa un ambiente $X$ e si limita la definizione agli $x\in X$. In quest'ultimo caso $\bigcap_{i\in\varnothing}A_i=X$; invece l'unione della famiglia vuota è sempre $\varnothing$.
\end{definizione}

% Integrazione M–Z L1 p. 2: proprietà dell'unione e dell'intersezione.
\begin{proposizione}
L'unione è il più piccolo insieme che contiene ogni $A_i$; l'intersezione è il più grande insieme contenuto in ogni $A_i$, nell'ambiente fissato.
\end{proposizione}
\begin{proof}
\textbf{Unione.} Ogni $x\in A_j$ appartiene all'unione scegliendo $i=j$, dunque $A_j\subseteq\bigcup_i A_i$. Se $A_i\subseteq B$ per ogni $i$, un elemento dell'unione sta in qualche $A_i$ e quindi in $B$: ne segue $\bigcup_i A_i\subseteq B$.

\textbf{Intersezione.} Ogni elemento dell'intersezione appartiene a ogni $A_j$. Se $B\subseteq A_i$ per ogni $i$, ciascun $x\in B$ appartiene a tutti gli $A_i$, dunque $B\subseteq\bigcap_i A_i$. Le convenzioni fissate danno le stesse proprietà anche per $I=\varnothing$.
\end{proof}

\begin{esercizio}[Associatività delle differenze]
Stabilire se la differenza e la differenza simmetrica sono associative.
\end{esercizio}
\begin{svolgimento}
\textbf{Differenza.} Per $A=B=C=\{0\}$ si ha $(A\setminus B)\setminus C=\varnothing$, mentre $A\setminus(B\setminus C)=\{0\}$. Quindi non è associativa. Non è neppure commutativa: $\{0\}\setminus\varnothing=\{0\}$ e $\varnothing\setminus\{0\}=\varnothing$.

\textbf{Differenza simmetrica.} Scomponiamo l'unione di tre insiemi secondo le appartenenze. Indichiamo con $A_0,B_0,C_0$ i punti appartenenti soltanto ad $A$, soltanto a $B$, soltanto a $C$; poniamo inoltre $T=A\cap B\cap C$. In formule,
\[
A_0=A\setminus(B\cup C),\quad B_0=B\setminus(A\cup C),\quad C_0=C\setminus(A\cup B).
\]
\passomath{(A\triangle B)\triangle C} I punti di $(A\triangle B)\setminus C$ stanno in uno solo fra $A,B$ e fuori da $C$: formano $A_0\cup B_0$. I punti di $C\setminus(A\triangle B)$ stanno in $C$ e, rispetto ad $A,B$, o in entrambi o in nessuno: formano $T\cup C_0$. Pertanto questo membro è $A_0\cup B_0\cup C_0\cup T$.

\passomath{A\triangle(B\triangle C)} I punti di $A\setminus(B\triangle C)$ stanno in $A$ e o in entrambi $B,C$ o in nessuno: formano $T\cup A_0$. I punti di $(B\triangle C)\setminus A$ formano $B_0\cup C_0$. Anche questo membro è dunque $A_0\cup B_0\cup C_0\cup T$. L'uguaglianza segue dall'associatività e dalla commutatività dell'unione.
\end{svolgimento}
% Ripresa effettiva A–L L3 pp. 1–2, confrontata con M–Z L2 pp. 5–6.

\begin{schema}
\begin{tikzpicture}[scale=.9]
  \draw (0,.45) circle (1.15);\draw (1.5,.45) circle (1.15);\draw (.75,-.85) circle (1.15);
  \node at (-.6,1.15) {$A$};\node at (2.1,1.15) {$B$};\node at (.75,-1.65) {$C$};
  \node at (-.45,.5) {$A_0$};\node at (1.95,.5) {$B_0$};\node at (.75,-1.3) {$C_0$};\node at (.75,-.05) {$T$};
  \node[align=center] at (.75,-2.5) {Le regioni $A_0$, $B_0$, $C_0$ e $T$\\formano l'unione comune ai due membri.};
\end{tikzpicture}
\end{schema}

\section{Prodotti cartesiani}
% A–L L1 pp. 3–4. Le potenze con insieme di indici sono rinviate alle funzioni.
\begin{definizione}
Il \textbf{prodotto cartesiano} è $A\times B=\{(a,b)\mid a\in A,\ b\in B\}$. Le coppie sono ordinate: $(a,b)=(a',b')$ se e solo se $a=a'$ e $b=b'$. Più in generale,
\[
A_1\times\cdots\times A_n=\{(a_1,\ldots,a_n)\mid a_i\in A_i\text{ per }i=1,\ldots,n\}.
\]
Se i fattori sono tutti $A$, si scrive $A^n$. In particolare $\R^2=\R\times\R$ è il piano delle coppie di numeri reali.
\end{definizione}

\begin{esempio}
Se $A=\{m,n,p\}$ e $B=\{9,8\}$, allora
\[
A\times B=\{(m,9),(m,8),(n,9),(n,8),(p,9),(p,8)\}.
\]
Il prodotto è vuoto se un fattore è vuoto; se entrambi sono non vuoti, scegliendo un elemento in ciascuno si ottiene una coppia nel prodotto. In generale $A\times B\ne B\times A$: scambiare le coordinate cambia le coppie, pur dando una corrispondenza uno a uno tra i due prodotti.
\end{esempio}

\section{Parti, complementi e identità}
% A–L L2 pp. 1–2; integrazioni logiche M–Z L2 pp. 2–6 e L3 p. 1.
\begin{definizione}
L'\textbf{insieme delle parti} di $X$ è $\mathcal P(X)=\{A\mid A\subseteq X\}$. Per $A\subseteq X$, il \textbf{complementare} nell'ambiente $X$ è $A^c=X\setminus A$, indicato anche con $C(A)$.
\end{definizione}

\begin{esempio}
Per $X=\{2,3,4\}$,
\[
\mathcal P(X)=\{\varnothing,\{2\},\{3\},\{4\},\{2,3\},\{2,4\},\{3,4\},X\}.
\]
Si distingue $A\subseteq X$ da $A\in\mathcal P(X)$: dicono la stessa cosa, ma la seconda considera $A$ come un elemento di un nuovo insieme.
\end{esempio}

\begin{schema}
\begin{tikzpicture}[scale=.75]
  \fill[coloreschema!16] (0,0) rectangle (4.5,2.4);\fill[white] (1.5,1.2) circle (.8);
  \draw (0,0) rectangle (4.5,2.4);\draw (1.5,1.2) circle (.8);
  \node at (1.5,1.2) {$A$};\node at (3.3,1.2) {$A^c$};\node[above right] at (4.5,2.2) {$X$};
\end{tikzpicture}
\end{schema}

\subsection{Logica delle appartenenze}
Fissato $x$, poniamo $\alpha=(x\in A)$ e $\beta=(x\in B)$. Le operazioni corrispondono ai connettivi logici secondo la tabella seguente; $\neg$ denota la negazione, $\lor$ l'«oppure» e $\land$ l'«e».
\[
\begin{array}{cc|cccc}
\alpha&\beta&x\in A\cup B&x\in A\cap B&x\in A\setminus B&x\in A\triangle B\\\hline
\mathrm V&\mathrm V&\mathrm V&\mathrm V&\mathrm F&\mathrm F\\
\mathrm V&\mathrm F&\mathrm V&\mathrm F&\mathrm V&\mathrm V\\
\mathrm F&\mathrm V&\mathrm V&\mathrm F&\mathrm F&\mathrm V\\
\mathrm F&\mathrm F&\mathrm F&\mathrm F&\mathrm F&\mathrm F
\end{array}
\]
Un'implicazione $\alpha\Rightarrow\beta$ è falsa soltanto quando $\alpha$ è vera e $\beta$ è falsa. Si verifica quindi, riga per riga, che
\[
(\alpha\Rightarrow\beta)\iff(\neg\alpha\lor\beta),\qquad
\neg(\alpha\Rightarrow\beta)\iff(\alpha\land\neg\beta).
\]
L'inclusione è un'affermazione universale: $A\subseteq B$ equivale a $\forall x\,(x\in A\Rightarrow x\in B)$. La sua negazione è $\exists x\,(x\in A\land x\notin B)$: un solo controesempio basta a negarla.

\begin{esercizio}[Negazione di un'implicazione]
La negazione di $\alpha\Rightarrow\beta$ coincide con un'implicazione il cui antecedente e conseguente siano scelti tra $\alpha,\beta,\neg\alpha,\neg\beta$?
\end{esercizio}
\begin{svolgimento}
La negazione cercata è vera in una sola delle quattro combinazioni di verità. Invece ogni implicazione del tipo indicato è vera almeno nelle due combinazioni in cui il suo antecedente è falso. Non possono quindi coincidere. La restrizione sulle formule ammesse è essenziale per la domanda.
\end{svolgimento}

\begin{esercizio}[Complementazione e leggi di De Morgan]
Verificare $(A^c)^c=A$ e dimostrare $(A\cup B)^c=A^c\cap B^c$, $(A\cap B)^c=A^c\cup B^c$.
\end{esercizio}
\begin{svolgimento}
Tutte le appartenenze sono considerate per $x\in X$. Non appartenere al complementare di $A$ equivale ad appartenere ad $A$, quindi $(A^c)^c=A$.

\passomath{(A\cup B)^c=A^c\cap B^c} Non appartenere all'unione significa non appartenere né ad $A$ né a $B$: $\neg(\alpha\lor\beta)\iff\neg\alpha\land\neg\beta$.

\passomath{(A\cap B)^c=A^c\cup B^c} Non appartenere all'intersezione significa che almeno una delle due appartenenze è falsa: $\neg(\alpha\land\beta)\iff\neg\alpha\lor\neg\beta$. Le equivalenze si possono anche controllare nelle quattro righe della tabella.
\end{svolgimento}

\begin{proposizione}[Distributività]
Per insiemi $A,B,C$ valgono
\[
\begin{aligned}
(A\cup B)\cap C&=(A\cap C)\cup(B\cap C),\\
(A\cap B)\cup C&=(A\cup C)\cap(B\cup C),\\
(A\triangle B)\cap C&=(A\cap C)\triangle(B\cap C).
\end{aligned}
\]
\end{proposizione}
\begin{proof}
Per la prima, $x$ deve appartenere a $C$ e ad almeno uno tra $A,B$: questo equivale a stare in $A\cap C$ oppure in $B\cap C$.

Per la seconda, se $x\in C$, appartiene a entrambi i membri. Se $x\notin C$, il primo richiede $x\in A\cap B$, e il secondo richiede simultaneamente $x\in A$ e $x\in B$: ancora la stessa condizione.

Per la terza, se $x\notin C$, non sta in nessuno dei due membri. Se $x\in C$, in entrambi i membri si richiede che stia in uno solo fra $A$ e $B$. I casi esauriscono le possibilità.
\end{proof}

\begin{esercizio}[Elementi neutri e annullamento]
In $\mathcal P(X)$ verificare $A\triangle\varnothing=A$ e $A\cap X=A$. Dato $A$, esiste $B$ tale che $A\triangle B=\varnothing$? Per quali $A$ esiste $B$ tale che $A\cap B=X$?
\end{esercizio}
\begin{svolgimento}
La differenza simmetrica con $\varnothing$ conserva esattamente gli elementi di $A$; l'intersezione con $X$ conserva tutti gli elementi di $A\subseteq X$. Questi sono gli elementi neutri per le due operazioni.

L'uguaglianza $A\triangle B=\varnothing$ dice che nessun elemento appartiene a uno solo tra $A,B$, cioè $A=B$. Quindi $B=A$ è sempre l'unica soluzione.

Se $A\cap B=X$, allora $X\subseteq A\subseteq X$ e $X\subseteq B\subseteq X$, dunque $A=B=X$. Viceversa questa scelta dà l'uguaglianza. Solo l'elemento $A=X$ ammette una soluzione per la seconda richiesta.
\end{svolgimento}

% !TeX root = ../main.tex
\chapter{Relazioni, equivalenze e quozienti}
\fonti{Lezioni 2, 3 e 4}

\section{Relazioni tra insiemi}
% A–L L2 pp. 3–4; M–Z L2 p. 7 e L3 pp. 1–3.
\begin{definizione}
Una \textbf{relazione} da $X$ a $Y$ è un sottoinsieme $\rho\subseteq X\times Y$. La scrittura $x\rho y$ significa $(x,y)\in\rho$. Se i due insiemi coincidono, si parla di relazione binaria \textbf{su} $X$.
\end{definizione}

\begin{esempio}[Città visitate]
Siano $X=\{P,V,R,M,T\}$ cinque città e $Y=\{\text{Anna},\text{Luca},\text{Lia}\}$. La relazione «la città è stata visitata dalla persona» può essere
\[
\begin{split}
\rho=\{&(P,\text{Anna}),(V,\text{Anna}),(R,\text{Anna}),\\
&(M,\text{Luca}),(T,\text{Luca}),(M,\text{Lia}),(R,\text{Lia})\}.
\end{split}
\]
La rappresentazione per tabella riporta le stesse coppie:
\[
\begin{array}{c|ccc}
&\text{Anna}&\text{Luca}&\text{Lia}\\\hline
P&\times&&\\V&\times&&\\R&\times&&\times\\M&&\times&\times\\T&&\times&
\end{array}
\]
Una città può essere in relazione con più persone e una persona con più città; una relazione non richiede alcuna unicità.
\end{esempio}

\begin{esercizio}[Grafici di relazioni]
Rappresentare $\rho=\{(x,y)\in\N\times\Z\mid x+y=1\}$, $\sigma=\{(x,y)\in\R^2\mid x^2+y^2=1\}$ e $\tau=\{(a,b)\in\N^*\times\N\mid b/a\in\N\}$.
\end{esercizio}
\begin{svolgimento}
Per $\rho$ si ottengono i punti $(n,1-n)$, $n\in\N$, a partire da $(0,1)$: è una parte discreta della retta $y=1-x$, non tutta la retta. Il grafico di $\sigma$ è la circonferenza di centro l'origine e raggio $1$. Per $\tau$, fissata un'ascissa intera $a>0$, le ordinate sono $b=ka$, $k\in\N$: tutti i multipli non negativi di $a$, compreso $0$.
\end{svolgimento}

\begin{schema}
\begin{tikzpicture}[scale=.65]
  \begin{scope}
    \draw[asse] (-.5,0)--(4,0) node[right] {$x$};\draw[asse] (0,-2.6)--(0,1.8) node[above] {$y$};
    \draw[densely dashed,black!45] (-.3,1.3)--(3.4,-2.4);
    \foreach \x in {0,...,3}{\fill[coloreschema] (\x,{1-\x}) circle (2pt);}
    \node at (1.6,-3.2) {$x+y=1$, $x\in\N$};
  \end{scope}
  \begin{scope}[shift={(7,0)}]
    \draw[asse] (-1.7,0)--(1.8,0) node[right] {$x$};\draw[asse] (0,-1.6)--(0,1.8) node[above] {$y$};
    \draw[coloreschema,thick] (0,0) circle (1);\node at (.9,.3) {$1$};
    \node at (0,-3.2) {$x^2+y^2=1$};
  \end{scope}
  \begin{scope}[shift={(11,0)}]
    \draw[asse] (-.3,0)--(4.5,0) node[right] {$a$};\draw[asse] (0,-.3)--(0,4.6) node[above] {$b$};
    \foreach \a in {1,...,4}{\foreach \b in {0,...,4}{\pgfmathtruncatemacro{\rem}{mod(\b,\a)}\ifnum\rem=0\fill[coloreschema] (\a,\b) circle (2pt);\fi}}
    \node at (2,-3.2) {$a\mid b$, porzione del grafico};
  \end{scope}
\end{tikzpicture}
\end{schema}

\section{Equivalenze e ordini}
\begin{definizione}
Una relazione $\rho$ su $X$ è \textbf{riflessiva} se $x\rho x$ per ogni $x\in X$; \textbf{simmetrica} se $x\rho y\Rightarrow y\rho x$ per ogni $x,y\in X$; \textbf{transitiva} se $x\rho y$ e $y\rho z$ implicano $x\rho z$ per ogni $x,y,z\in X$. Una \textbf{relazione di equivalenza} soddisfa tutte e tre le proprietà.
\end{definizione}

\begin{esempio}[Uguaglianza]
La diagonale $\Delta_X=\{(x,x)\mid x\in X\}$ definisce l'uguaglianza: $x=x$; da $x=y$ segue $y=x$; da $x=y$ e $y=z$ segue $x=z$. Sono precisamente le tre proprietà richieste.
\end{esempio}

\begin{esercizio}[Tre equivalenze]
Verificare che sono equivalenze la relazione $(a,b)\sim(c,d)\iff ad=bc$ su $F=\Z\times\Z^*$, la similitudine tra triangoli non degeneri del piano e il parallelismo tra rette, ammettendo rette coincidenti.
\end{esercizio}
\begin{svolgimento}
\textbf{Coppie intere.} La riflessività è $ab=ba$. La simmetria segue da $ad=bc\Rightarrow cb=da$. Per la transitività, se $(a,b)\sim(c,d)$ e $(c,d)\sim(e,f)$, allora $ad=bc$ e $cf=de$. Moltiplicando si ottiene $adf=bcf=bde$; poiché $d\ne0$, si cancella $d$ e si ricava $af=be$, cioè $(a,b)\sim(e,f)$.

\textbf{Triangoli.} Per il criterio degli angoli, due triangoli sono simili se i rispettivi angoli si possono mettere in corrispondenza con ampiezze uguali. Ogni triangolo ha gli stessi angoli di se stesso; una corrispondenza di angoli uguali vale anche in senso contrario; se gli angoli del primo coincidono con quelli del secondo e quelli del secondo con quelli del terzo, coincidono anche quelli del primo e del terzo. Si hanno riflessività, simmetria e transitività.

\textbf{Rette.} Essere parallele, ammettendo la coincidenza, significa avere la stessa direzione. L'uguaglianza delle direzioni è riflessiva, simmetrica e transitiva. Escludendo la coincidenza si perderebbe la riflessività.
\end{svolgimento}

\begin{esercizio}[Riconoscere le proprietà]
Stabilire quali delle seguenti relazioni sono equivalenze: $a\rho b\iff b=a^2$ su $\R$; la perpendicolarità tra rette; $a\sigma b\iff a\le b$ su $\N$; $a\tau b\iff |a|=|b|$ su $\R$; $x\sim_f y\iff f(x)=f(y)$, dove a ogni elemento $x$ è associato un unico valore $f(x)$.
\end{esercizio}
\begin{svolgimento}
\textbf{Quadrato.} Non è riflessiva ($2\not\rho2$), né simmetrica ($2\rho4$, ma $4\not\rho2$), né transitiva ($2\rho4$ e $4\rho16$, ma $2\not\rho16$).

\textbf{Perpendicolarità.} È simmetrica, ma una retta non è perpendicolare a se stessa. Non è transitiva: da $r\perp s$ e $s\perp r$ non segue $r\perp r$.

\textbf{Ordine usuale.} È riflessivo e transitivo, ma $0\le1$ e $1\nleq0$, quindi non simmetrico.

\textbf{Valore assoluto e uguaglianza dei valori.} Sono entrambe equivalenze: si applicano le tre proprietà dell'uguaglianza, rispettivamente ai numeri $|a|$ e ai valori $f(x)$. In particolare $f(x)=f(y)=f(z)$ implica $f(x)=f(z)$, senza richiedere che $x=y=z$.
\end{svolgimento}

\begin{definizione}
Una relazione è \textbf{antisimmetrica} se $x\rho y$ e $y\rho x$ implicano $x=y$. Un \textbf{ordine parziale} è una relazione riflessiva, antisimmetrica e transitiva. È un \textbf{ordine totale} se ogni coppia è confrontabile: per ogni $x,y$ vale $x\rho y$ oppure $y\rho x$.
\end{definizione}

\begin{esercizio}[Esempi di ordini]
Verificare che l'inclusione è un ordine parziale su $\mathcal P(X)$, la divisibilità è un ordine parziale su $\N^*$ e $\le$ è un ordine totale su $\R$. Che cosa cambia per la divisibilità su $\Z^*$?
\end{esercizio}
\begin{svolgimento}
\textbf{Inclusione.} $A\subseteq A$; $A\subseteq B\subseteq A$ implica $A=B$; $A\subseteq B\subseteq C$ implica $A\subseteq C$. Le tre proprietà seguono direttamente dall'appartenenza.

\textbf{Divisibilità.} Scriviamo $a\mid b$ se $b=ak$ per qualche intero $k$. Si ha $a=a\cdot1$. Se $b=ak$ e $c=bh$, allora $c=a(kh)$, quindi vale la transitività. Per $a,b>0$, da $a\mid b$ segue $a\le b$; se anche $b\mid a$, si ha $a=b$. Non è totale: $2$ e $3$ non sono confrontabili. Su $\Z^*$ viene meno l'antisimmetria: $1\mid-1$ e $-1\mid1$, ma $1\ne-1$.

\textbf{Ordine reale.} Le proprietà usuali di $\le$ danno riflessività, antisimmetria, transitività e confrontabilità di ogni coppia, quindi un ordine totale.
\end{svolgimento}

\begin{esercizio}[Confrontabilità]
Per $X\ne\varnothing$, l'inclusione su $\mathcal P(X)$ è sempre totale? Su $S=\{1,2,-2,6,-6,18\}$, con la divisibilità intera, ogni coppia è confrontabile? Vale l'antisimmetria?
\end{esercizio}
\begin{svolgimento}
Se $X$ contiene $x\ne y$, i singoletti $\{x\}$ e $\{y\}$ non sono confrontabili. Se $X$ ha un solo elemento, invece, $\mathcal P(X)=\{\varnothing,X\}$ è una catena.

In $S$ i valori assoluti sono ordinati dalla divisibilità come $1\mid2\mid6\mid18$. Il segno non modifica la divisibilità tra interi, perciò ogni coppia è confrontabile. Ma $2\mid-2$ e $-2\mid2$ con $2\ne-2$: confrontabilità e antisimmetria sono condizioni diverse, e qui non si ha un ordine totale.
\end{svolgimento}

\section{Partizioni, classi e quozienti}
% A–L L3 pp. 1, 3–7; M–Z L3 pp. 5–9. Conservata la precedenza A–L della partizione.
\begin{definizione}
Una \textbf{partizione} di $X$ è un insieme $\mathcal A$ di sottoinsiemi non vuoti di $X$, detti \textbf{blocchi}, la cui unione è $X$ e tali che due blocchi distinti siano disgiunti. Equivalentemente,
\[
\bigcup_{A\in\mathcal A}A=X,\qquad
A\cap B\ne\varnothing\ \Longleftrightarrow\ A=B
\quad(A,B\in\mathcal A).
\]
Nella seconda formulazione, prendendo $A=B$, si ricava anche che ogni blocco è non vuoto. Ogni elemento di $X$ appartiene a un unico blocco. Per $X=\varnothing$ la partizione è la famiglia vuota.
\end{definizione}

\begin{esempio}[Scomposizione di un'unione]
I tre pezzi $A\setminus B$, $A\cap B$, $B\setminus A$ ricoprono $A\cup B$ e sono a due a due disgiunti. Eliminando i pezzi vuoti formano una partizione. Per tre insiemi si ottengono analogamente le sette regioni determinate dalle appartenenze, alcune eventualmente vuote. La decomposizione usata per la differenza simmetrica sceglie quattro di queste regioni.
\end{esempio}

\begin{definizione}
Per un'equivalenza $\sim$ su $X$, la \textbf{classe di equivalenza} di $x\in X$ è $[x]=\{y\in X\mid y\sim x\}$. Ogni suo elemento è un \textbf{rappresentante} della classe. L'\textbf{insieme quoziente} è $X/{\sim}=\{[x]\mid x\in X\}\subseteq\mathcal P(X)$.
\end{definizione}

\begin{proposizione}\label{prop:classi}
Per $x,y\in X$, $x\sim y$ se e solo se $[x]=[y]$.
\end{proposizione}
\begin{proof}
\passomath{x\sim y\Rightarrow[x]\subseteq[y]} Se $z\in[x]$, allora $z\sim x$ e $x\sim y$; per transitività $z\sim y$, quindi $z\in[y]$.

\passomath{x\sim y\Rightarrow[y]\subseteq[x]} Per simmetria $y\sim x$. Se $z\in[y]$, da $z\sim y\sim x$ segue $z\in[x]$. Le due inclusioni danno l'uguaglianza.

\passomath{[x]=[y]\Rightarrow x\sim y} Per riflessività $x\in[x]$; dall'uguaglianza segue $x\in[y]$, che significa precisamente $x\sim y$.
\end{proof}

\begin{teorema}\label{teo:partizioni}
Le classi di un'equivalenza formano una partizione di $X$. Viceversa, se $\mathcal A$ è una partizione, la relazione $x\sim_{\mathcal A}y$ definita dall'appartenenza a uno stesso blocco è un'equivalenza, le cui classi sono i blocchi di $\mathcal A$.
\end{teorema}
\begin{proof}
\textbf{Dall'equivalenza alla partizione.} Ogni classe contiene il proprio rappresentante, dunque non è vuota. Ogni $x\in X$ appartiene a $[x]$, perciò l'unione delle classi contiene $X$; l'altra inclusione vale perché ogni classe è contenuta in $X$. Se $z\in[x]\cap[y]$, si ha $z\sim x$ e $z\sim y$. Per simmetria e transitività segue $x\sim z\sim y$, quindi $[x]=[y]$. Due classi distinte sono pertanto disgiunte.

\textbf{Dalla partizione all'equivalenza.} Ogni $x$ appartiene a un blocco, dunque $x\sim_{\mathcal A}x$. Se $x,y$ stanno nello stesso blocco, vi stanno anche $y,x$, quindi la relazione è simmetrica. Se $x,y\in A$ e $y,z\in B$, i blocchi si incontrano in $y$ e perciò $A=B$; allora $x,z$ stanno nello stesso blocco, ottenendo la transitività.

\textbf{Identificazione delle classi.} Sia $x\in A\in\mathcal A$. Ogni $y\in A$ è equivalente a $x$, quindi $A\subseteq[x]$. Viceversa, se $y\sim_{\mathcal A}x$, un blocco $B$ contiene entrambi; poiché $x\in A\cap B$, si ha $A=B$ e $y\in A$. Dunque $[x]=A$.
\end{proof}

\begin{schema}
\begin{tikzpicture}
  \draw[rounded corners] (0,0) rectangle (6,2.4);
  \draw (2,0)--(2,2.4);\draw (4,0)--(4,2.4);
  \node at (1,1.9) {$[x]$};\node at (3,1.9) {$[y]$};\node at (5,1.9) {$[z]$};
  \foreach \x/\y in {.5/.7,1.4/1,2.6/.6,3.4/1.1,4.6/.7,5.3/1.1}{\fill[coloreschema] (\x,\y) circle (1.5pt);}
  \node[right] at (6.2,1.2) {$X$};
\end{tikzpicture}
\end{schema}

\begin{esercizio}[Le costruzioni inverse]
Verificare che le due costruzioni del teorema stabiliscono una corrispondenza uno a uno tra equivalenze su $X$ e partizioni di $X$.
\end{esercizio}
\begin{svolgimento}
Partendo da una partizione, il teorema mostra che le classi della relazione costruita sono esattamente i blocchi di partenza: si recupera la partizione.

Partendo da $\sim$, due elementi $x,y$ appartengono a una stessa classe se e solo se $[x]=[y]$, cioè, per la proposizione~\ref{prop:classi}, se e solo se $x\sim y$. Si recupera anche la relazione iniziale. Ciascuna costruzione annulla dunque l'altra.
\end{svolgimento}

\begin{esercizio}[Descrivere i quozienti]
Descrivere le classi e i quozienti degli esempi di equivalenza precedenti, compresa la relazione $x^2=y^2$ su $\R$.
\end{esercizio}
\begin{svolgimento}
Per l'uguaglianza, le classi sono i singoletti $\{x\}$. Per $|x|=|y|$, sono $\{0\}$ e $\{t,-t\}$ con $t>0$; lo stesso vale per $x^2=y^2$, perché $x^2-y^2=(x-y)(x+y)$. Per esempio $[5]=[-5]=\{5,-5\}$.

Per il parallelismo, una classe comprende tutte le rette di una stessa direzione. Per la similitudine, comprende tutti i triangoli di una stessa forma: la terna degli angoli, senza fissarne l'ordine, determina la classe. Per l'uguaglianza $f(x)=f(y)$, una classe raccoglie tutti gli elementi ai quali è associato lo stesso valore. Il quoziente delle coppie intere per prodotti incrociati è descritto di seguito.
\end{svolgimento}

\section{I razionali e le operazioni sul quoziente}
% A–L L3 p. 8 e L4 pp. 1–3; M–Z L4 pp. 1–4.
\begin{definizione}
Posto $F=\Z\times\Z^*$ e $(a,b)\sim(c,d)\iff ad=bc$, si definisce $\Q=F/{\sim}$. La notazione $a/b$ indica la classe $[(a,b)]$, non la singola coppia. Per esempio
\[
\frac{-2}{4}=\frac{-1}{2}=\frac{1}{-2},\qquad
\frac12=\frac24=\frac{-2}{-4}=\frac36.
\]
\end{definizione}

Una regola sui rappresentanti dà un'operazione sulle classi solo se il risultato è indipendente dai rappresentanti scelti.

\begin{esempio}[Una regola incompatibile]
La formula $a/b\star c/d=(a+c)/(bd)$ non definisce un'operazione su $\Q$: infatti $1/2=2/4$, ma
\[
\frac12\star\frac23=\frac36=\frac12,\qquad
\frac24\star\frac23=\frac4{12}=\frac13.
\]
La stessa coppia di razionali produrrebbe due risultati diversi.
\end{esempio}

% Mantiene insieme enunciato e sviluppo algebrico sul raccordo finale.
\Needspace{30\baselineskip}
\begin{proposizione}[Somma e prodotto]
Le formule
\[
\frac ab+\frac cd=\frac{ad+bc}{bd},\qquad
\frac ab\cdot\frac cd=\frac{ac}{bd}
\]
definiscono operazioni su $\Q$.
\end{proposizione}
\begin{proof}
Siano $ab'=a'b$ e $cd'=c'd$, con $b,b',d,d'\ne0$. I denominatori $bd$ e $b'd'$ sono non nulli, dunque le coppie risultanti sono ammesse.

\textbf{Somma.} Per l'equivalenza delle coppie occorre $(ad+bc)b'd'=(a'd'+b'c')bd$. Le ipotesi danno
\[
\begin{aligned}
(ad+bc)b'd'&=(ab')dd'+bb'(cd')\\
&=(a'b)dd'+bb'(c'd)\\
&=(a'd'+b'c')bd.
\end{aligned}
\]
\textbf{Prodotto.} Analogamente $(ac)b'd'=(ab')(cd')=(a'b)(c'd)=(a'c')bd$. Entrambi i risultati appartengono quindi alla stessa classe dopo il cambio di rappresentanti.
\end{proof}

\begin{proposizione}[Criterio di compatibilità]
Sia $\star$ un'operazione su $X$, cioè una regola che a ogni coppia di elementi di $X$ associa un elemento di $X$. La formula $[x]\mathbin{\overline\star}[y]=[x\star y]$ definisce un'operazione su $X/{\sim}$ se e solo se
\[
x\sim x',\ y\sim y'\quad\Longrightarrow\quad x\star y\sim x'\star y'
\]
per ogni $x,x',y,y'\in X$.
\end{proposizione}
\begin{proof}
\passomath{\Rightarrow} Se $x\sim x'$ e $y\sim y'$, le classi in ingresso coincidono. Un'operazione ben definita deve dare la stessa classe in uscita, ossia $[x\star y]=[x'\star y']$; per la proposizione~\ref{prop:classi} i rappresentanti sono equivalenti.

\passomath{\Leftarrow} Due scelte dei rappresentanti di $[x]$ e $[y]$ differiscono per le equivalenze $x\sim x'$ e $y\sim y'$. La condizione assicura che i risultati abbiano la stessa classe. La regola è quindi indipendente da entrambe le scelte.
\end{proof}

% !TeX root = ../main.tex
\chapter{Funzioni}
\fonti{Lezioni 2, 4, 5, 6 e 7}

\section{Dominio, codominio e grafico}
% A–L L4 pp. 3–5; M–Z L4 pp. 4–7. L'ordine resta A–L.
\begin{definizione}
Una \textbf{funzione} $f:X\to Y$ consiste di un dominio $X$, un codominio $Y$ e una regola che associa a ogni $x\in X$ un unico elemento $f(x)\in Y$. L'\textbf{immagine} e il \textbf{grafico} sono
\[
\operatorname{Im}(f)=\{f(x)\mid x\in X\}\subseteq Y,\qquad
\operatorname{Gr}(f)=\{(x,f(x))\mid x\in X\}\subseteq X\times Y.
\]
Due funzioni sono uguali quando coincidono dominio, codominio e valori in ciascun punto del dominio.
\end{definizione}

\begin{esempio}[Regole sui rappresentanti]
La formula $f:\Q\to\Q$, $a/b\mapsto a+b$, non definisce una funzione: la stessa classe $1/2=2/4$ darebbe $3$ e $6$. Anche $(a/b,c/d)\mapsto(a+c)/(bd)$ fallisce per il controesempio del capitolo precedente. Una formula non basta: deve assegnare un unico risultato all'elemento del dominio, indipendentemente da come lo si scrive.
\end{esempio}

\begin{esempio}[Sei funzioni]
Useremo le seguenti applicazioni:
\begin{enumerate}[label=\alph*)]
\item $c:\C\to\C$, $c(z)=0$;
\item $\operatorname{id}_{\R}:\R\to\R$, $x\mapsto x$;
\item $s:\R_{\ge0}\to\R_{\ge0}$, $s(x)=x^2$;
\item $a:T\to\R$, che associa a un triangolo non degenere del piano la sua area;
\item $d:\N\times\N\to\Z$, $d(m,n)=m-n$;
\item $J:\N^*\to\mathcal P(P)$, $J(n)=\{p\in P\mid p\mid n\}$, dove $P$ è l'insieme dei primi positivi e $p\mid n$ significa che $n$ è un multiplo intero di $p$.
\end{enumerate}
Per l'ultima, $J(1)=\varnothing$, $J(4)=\{2\}$, $J(12)=\{2,3\}$, $J(30)=\{2,3,5\}$, $J(7)=\{7\}$. In questo esempio si usano le proprietà aritmetiche usuali dei primi come fatti assunti: la loro infinità e la fattorizzazione saranno dimostrate nel capitolo~\ref{cap:interi}.
\end{esempio}

Il grafico è una relazione particolare: per ogni $x\in X$ deve esistere esattamente un $y\in Y$ con $(x,y)$ nel grafico. La condizione riguarda tutti i punti del dominio, ma non impone che ogni punto del codominio sia raggiunto. Nel piano, ogni retta verticale relativa a un'ascissa del dominio deve incontrare il grafico una sola volta.

\begin{schema}
\begin{tikzpicture}[scale=.72]
\begin{scope}
\draw[asse] (-1.8,0)--(1.8,0) node[right] {$x$};\draw[asse] (0,-.3)--(0,2.6) node[above] {$y$};
\draw[coloreschema,thick,domain=-1.45:1.45,samples=50] plot (\x,{\x*\x});\draw[dashed] (1,-.15)--(1,2.3);\fill (1,1) circle (2pt);\node at (0,-.8) {Un unico punto per ascissa};
\end{scope}
\begin{scope}[shift={(7,1)}]
\draw[asse] (-1.8,0)--(1.8,0) node[right] {$x$};\draw[asse] (0,-1.4)--(0,1.7) node[above] {$y$};
\draw[coloreschema,thick] (0,0) circle (1);\draw[dashed] (.5,-1.3)--(.5,1.3);\fill (.5,.866) circle (2pt);\fill (.5,-.866) circle (2pt);\node at (0,-1.8) {La circonferenza non è un grafico};
\end{scope}
\end{tikzpicture}
\end{schema}

\begin{definizione}
Una funzione $f:X\to Y$ è \textbf{iniettiva} se, per ogni $x,x'\in X$, $f(x)=f(x')$ implica $x=x'$. È \textbf{suriettiva} se per ogni $y\in Y$ esiste $x\in X$ con $f(x)=y$, ossia $\operatorname{Im}(f)=Y$.
\end{definizione}

\begin{esercizio}[Immagine, iniettività e suriettività]
Determinare l'immagine e discutere iniettività e suriettività delle sei funzioni precedenti. Per il quadrato considerare anche i domini $\R$ e $\R_{\le0}$ e il codominio $\R$.
\end{esercizio}
\begin{svolgimento}
\textbf{Costante.} L'immagine è $\{0\}$, diversa da $\C$, e $c(0)=c(1)$: non è né iniettiva né suriettiva.

\textbf{Identità.} $\operatorname{id}_{\R}(x)=\operatorname{id}_{\R}(x')$ è già $x=x'$, e ogni $y$ è immagine di se stesso. È iniettiva e suriettiva.

\textbf{Quadrato.} Per $x,x'\ge0$, da $x^2=(x')^2$ segue $(x-x')(x+x')=0$. Se $x+x'=0$, entrambi sono $0$; altrimenti si cancella il secondo fattore e si ha $x=x'$. Ogni $y\ge0$ è immagine di $\sqrt y$, perciò $s$ è iniettiva e suriettiva. Anche $\R_{\le0}\to\R_{\ge0}$, $x\mapsto x^2$, è iniettiva: due numeri non positivi con lo stesso quadrato coincidono; è suriettiva perché $y$ è immagine di $-\sqrt y$. Su tutto $\R$ non è iniettiva, dato che $1$ e $-1$ hanno la stessa immagine; con codominio $\R$ non è suriettiva, perché l'immagine rimane $[0,+\infty)$.

\textbf{Area.} Due triangoli distinti ottenuti per traslazione hanno la stessa area, quindi non c'è iniettività. Ogni $t>0$ è area di un triangolo di base $2$ e altezza $t$, mentre un triangolo non degenere ha area positiva. L'immagine è $(0,+\infty)$, diversa dal codominio $\R$.

\textbf{Differenza.} $d(0,0)=d(1,1)$, quindi non è iniettiva. Ogni $z\ge0$ è $d(z,0)$ e ogni $z<0$ è $d(0,-z)$: l'immagine è $\Z$ e la funzione è suriettiva.

\textbf{Divisori primi.} $J(2)=J(4)$, quindi non è iniettiva. Per $n>0$, i primi che dividono $n$ sono compresi tra $2$ e $n$, dunque $J(n)$ è finito. Viceversa un insieme finito $\{p_1,\ldots,p_r\}$ di primi è immagine del prodotto $p_1\cdots p_r$; per l'insieme vuoto si usa $n=1$. Quindi l'immagine è l'insieme dei sottoinsiemi finiti di $P$, non tutto $\mathcal P(P)$, poiché $P$ è infinito. Non è suriettiva.
\end{svolgimento}

\begin{esercizio}[Funzione a tratti sugli interi]
Studiare $f:\Z\to\Z$ definita da
\[
f(n)=\begin{cases}2n-3&n\text{ pari},\\n+1&n\text{ dispari}.\end{cases}
\]
\end{esercizio}
\begin{svolgimento}
% A–L L4 p. 5; svolgimento M–Z L7 pp. 1–2, ripresa effettiva.
\textbf{Iniettività.} Un intero pari $n=2k$ ha immagine $4k-3$, dispari; un intero dispari $n=2k+1$ ha immagine $2k+2$, pari. Se due elementi hanno immagini uguali, devono quindi appartenere allo stesso ramo. Nel ramo pari, $2n-3=2n'-3$ implica $n=n'$; in quello dispari, $n+1=n'+1$ implica ancora $n=n'$. Anche i casi misti sono così esclusi.

\textbf{Immagine.} Il ramo dispari percorre tutti gli interi pari: dato $y$ pari, $n=y-1$ è dispari e $f(n)=y$. Il ramo pari percorre esattamente gli interi $4k-3$, $k\in\Z$. Non raggiunge $3$, perché $4k-3=3$ richiederebbe $4k=6$. Dunque $f$ è iniettiva ma non suriettiva.
\end{svolgimento}

\begin{esercizio}[Funzione del piano]
Studiare $F:\R^2\to\R^2$, $F(x,y)=(y^3,x(y^2+1))$.
\end{esercizio}
\begin{svolgimento}
% A–L L4 p. 6; prova completa M–Z L5 pp. 2–3.
\textbf{Iniettività.} Se $F(x,y)=F(x',y')$, l'uguaglianza della prima coordinata dà $y^3=(y')^3$, quindi $y=y'$. La seconda diventa $x(y^2+1)=x'(y^2+1)$. Poiché $y^2+1>0$, si cancella e si ricava $x=x'$.

\textbf{Suriettività.} Fissato un punto arbitrario $(u,v)$ del codominio, le equazioni da risolvere sono $y^3=u$ e $x(y^2+1)=v$. Poniamo
\[
y=\sqrt[3]u,\qquad x=\frac{v}{1+(\sqrt[3]u)^2}.
\]
Sono numeri reali ben definiti, perché il denominatore è positivo. Sostituendo si ottiene $F(x,y)=(u,v)$: ogni punto ha una preimmagine.
\end{svolgimento}

\section{Fibre e proiezione canonica}
% A–L L4 pp. 6–7 e L5 p. 1; ritorno J su N in A–L L6 pp. 1–2.
\begin{definizione}
La \textbf{fibra} di $f:X\to Y$ sopra $y\in Y$ è $\{x\in X\mid f(x)=y\}$. L'equivalenza $x\sim_f x'\iff f(x)=f(x')$ ha come classe di $x$ la fibra sopra $f(x)$:
\[
[x]_{\sim_f}=\{x'\in X\mid f(x')=f(x)\}.
\]
Le fibre non vuote formano una partizione di $X$. Quelle sopra i punti di $Y\setminus\operatorname{Im}(f)$ sono vuote e non sono classi.
\end{definizione}

\begin{schema}
\begin{tikzpicture}[>=Stealth]
  \draw (0,0) ellipse (1.4 and 1.15);\draw (5,0) ellipse (1.1 and 1.15);
  \node at (-1.3,1.2) {$X$};\node at (6,1.2) {$Y$};
  \draw[dashed,coloreschema] (-.2,.25) ellipse (.8 and .65);
  \node (a) at (-.35,.6) {$x_1$};\node (b) at (-.15,-.1) {$x_2$};\node (c) at (0,-.8) {$x_3$};
  \node (u) at (4.8,.6) {$y_0$};\node (v) at (5,-.15) {$y_1$};\node at (5,-.85) {$y_2$};
  \draw[->,coloreschema] (a)--(u);\draw[->,coloreschema] (b)--(u);\draw[->,coloreschema] (c)--(v);
  \node at (2.5,-1.65) {Una fibra può contenere più punti; un valore può non avere preimmagini.};
\end{tikzpicture}
\end{schema}

\begin{proposizione}
La funzione $f$ è iniettiva se e solo se $\sim_f$ è l'uguaglianza su $X$.
\end{proposizione}
\begin{proof}
Se $f$ è iniettiva, $x\sim_f x'$ implica $x=x'$; il viceversa vale sempre, quindi le due relazioni coincidono. Se invece coincidono, da $f(x)=f(x')$ segue $x\sim_f x'$, cioè $x=x'$: è l'iniettività.
\end{proof}

\begin{esempio}
Per $x\mapsto x^2$ su $\R$, le fibre non vuote sono $\{0\}$ e $\{-\sqrt y,\sqrt y\}$ per $y>0$. In particolare $[2]=[-2]=\{2,-2\}$. Sopra $-1$ la fibra è vuota.
\end{esempio}

\begin{esercizio}[Le fibre dei divisori primi]
Estendere $J$ a $\N$ con la stessa formula. Determinarne l'immagine, l'iniettività, la suriettività e le classi di $0,1,4,5,6,12,168$.
\end{esercizio}
\begin{svolgimento}
Ogni primo divide $0$, quindi $J(0)=P$. Le altre immagini sono quelle già trovate: tutti e soli i sottoinsiemi finiti di $P$. Pertanto $\operatorname{Im}(J)=\{A\subseteq P\mid A\text{ finito}\}\cup\{P\}$. Rimane non iniettiva perché $J(2)=J(4)$, e non suriettiva perché $P\setminus\{2\}$ è infinito e diverso da $P$.

Usando la fattorizzazione in primi, qui ancora assunta, si ottengono
\[
\begin{gathered}
[0]=\{0\},\qquad [1]=\{1\},\qquad [4]=\{2^k\mid k\ge1\},\qquad
[5]=\{5^k\mid k\ge1\},\\
[6]=[12]=\{2^a3^b\mid a,b\ge1\},\\
[168]=\{2^a3^b7^c\mid a,b,c\ge1\}.
\end{gathered}
\]
Infatti $168=2^3\cdot3\cdot7$. In ogni classe devono comparire tutti i primi indicati, dunque gli esponenti sono positivi, e non possono comparire primi diversi. La classe di $0$ è un singoletto perché ogni $n>0$ ha solo un numero finito di divisori primi.
\end{svolgimento}

\begin{definizione}
Una funzione è \textbf{biettiva}, o una \textbf{biezione}, quando è sia iniettiva sia suriettiva.
\end{definizione}

\begin{proposizione}[Proiezione canonica]
Se $\sim$ è un'equivalenza su $X$, la funzione $\pi:X\to X/{\sim}$, $x\mapsto[x]$, è suriettiva. Le sue fibre sono le classi dell'equivalenza iniziale; in particolare $\sim_\pi$ coincide con $\sim$.
\end{proposizione}
\begin{proof}
Ogni elemento del quoziente è una classe $[x]$, dunque è immagine di $x$. Inoltre $\pi(x')=\pi(x)$ equivale a $[x']=[x]$, cioè a $x'\sim x$ per la proposizione~\ref{prop:classi}. La fibra sopra $[x]$ è proprio il sottoinsieme $[x]\subseteq X$, e l'equivalenza indotta è quella di partenza.
\end{proof}

\section{Immagine diretta e controimmagine}
% A–L L5 p. 2; M–Z L5 pp. 6–8.
\begin{definizione}
A $f:X\to Y$ si associano le applicazioni
\[
\begin{aligned}
f_*:\mathcal P(X)&\to\mathcal P(Y),& f_*(A)&=\{f(a)\mid a\in A\},\\
f^*:\mathcal P(Y)&\to\mathcal P(X),& f^*(B)&=\{x\in X\mid f(x)\in B\},
\end{aligned}
\]
dette \textbf{immagine diretta} e \textbf{controimmagine}. La scrittura $f(A)$ denota spesso $f_*(A)$; la fibra sopra $y$ è $f^*(\{y\})$. La controimmagine esiste per qualunque funzione, senza supporre invertibilità.
\end{definizione}

\begin{esempio}[Il quadrato]
Per $f:\R\to\R$, $f(x)=x^2$, si ha
\[
f_*([-1,2])=f_*([-2,0])=[0,4],\qquad
f^*([0,4])=f^*((-\infty,4])=[-2,2].
\]
I quadrati nei due intervalli percorrono $[0,4]$; la condizione $x^2\le4$ equivale a $-2\le x\le2$, mentre $x^2\ge0$ è automatica. Le varianti $[-2,1]$ per l'immagine e $[-2,4]$ per la controimmagine danno gli stessi risultati. Invece $f^*([-3,-1])=\varnothing$, perché un quadrato non è negativo.

Per $A=[-1,2]$ si ottiene $f^*(f_*(A))=[-2,2]\ne A$: la controimmagine recupera tutte le fibre che incontrano $A$, anche dove esse escono da $A$.
\end{esempio}

\begin{schema}
\begin{tikzpicture}[scale=.85]
  \draw[asse] (-2.8,0)--(2.8,0) node[right] {$x$};\draw[asse] (0,-.3)--(0,4.9) node[above] {$y$};
  \draw[coloreschema,thick,domain=-2.1:2.1,samples=60] plot (\x,{\x*\x});
  \draw[densely dashed] (-2,4)--(2,4);\draw[densely dashed] (-2,0)--(-2,4);\draw[densely dashed] (2,0)--(2,4);
  \node[below] at (-2,0) {$-2$};\node[below] at (-1,0) {$-1$};\node[below] at (2,0) {$2$};\node[above left] at (0,4) {$4$};
  \draw[coloreschema,line width=2pt] (-1,-.35)--(2,-.35);\node[right] at (2.2,-.35) {$A$};
\end{tikzpicture}
\end{schema}

\begin{esercizio}[Confronto fra immagini e controimmagini]
Dimostrare, per ogni $A\subseteq X$ e $B\subseteq Y$,
\[
A\subseteq f^*(f_*(A)),\qquad f_*(f^*(B))\subseteq B.
\]
Quando la prima è un'uguaglianza per ogni $A$? Quando la seconda è un'uguaglianza per ogni $B$?
\end{esercizio}
\begin{svolgimento}
\textbf{Inclusioni.} Se $x\in A$, allora $f(x)\in f_*(A)$, dunque $x\in f^*(f_*(A))$. Per l'altra inclusione si ha più precisamente
\[
y\in f_*(f^*(B))\iff \exists x\in X:\ y=f(x)\in B
\iff y\in B\cap\operatorname{Im}(f).
\]
Quindi $f_*(f^*(B))=B\cap\operatorname{Im}(f)\subseteq B$.

\textbf{Prima uguaglianza.} Se $f$ è iniettiva e $x\in f^*(f_*(A))$, esiste $a\in A$ con $f(x)=f(a)$, quindi $x=a\in A$. Vale anche l'inclusione inversa, dunque l'uguaglianza. Viceversa, supponiamo l'uguaglianza per ogni $A$ e siano $f(x)=f(x')$. Per $A=\{x\}$, $x'\in f^*(f_*(A))=A$, quindi $x'=x$. La condizione equivale all'iniettività.

\textbf{Seconda uguaglianza.} Se $f$ è suriettiva, $\operatorname{Im}(f)=Y$ e la formula precedente dà l'uguaglianza per ogni $B$. Viceversa basta applicarla a $B=Y$: si ricava $f_*(X)=Y$. La condizione equivale alla suriettività.
\end{svolgimento}

\section{Composizione}
% A–L L5 pp. 3–6; M–Z L6 pp. 1–4.
\begin{definizione}
Se $f:X\to Y$ e $g:Y\to Z$, si definisce la \textbf{composizione} $g\circ f:X\to Z$ mediante $(g\circ f)(x)=g(f(x))$: si applica prima $f$, poi $g$.
\end{definizione}

Il codominio della prima funzione è il dominio della seconda, quindi il valore $f(x)$ è sempre un argomento ammesso per $g$. Per esempio, ponendo $f:\{1,2,3\}\to\N$ con $f(1)=0$, $f(2)=f(3)=1$, e $g:\N\to\R$, $g(n)=n$, la composta ha gli stessi valori, ma codominio $\R$.

\begin{proposizione}[Associatività]
Per $f:X\to Y$, $g:Y\to Z$ e $h:Z\to T$, si ha $h\circ(g\circ f)=(h\circ g)\circ f$.
\end{proposizione}
\begin{proof}
Le due funzioni hanno dominio $X$ e codominio $T$. Per ogni $x\in X$,
\[
(h\circ(g\circ f))(x)=h(g(f(x)))=((h\circ g)\circ f)(x).
\]
L'uguaglianza in ogni punto conclude la prova.
\end{proof}

\begin{proposizione}
Per $f:X\to Y$ e $g:Y\to Z$ valgono
\[
(g\circ f)_*=g_*\circ f_*:\mathcal P(X)\to\mathcal P(Z),\qquad
(g\circ f)^*=f^*\circ g^*:\mathcal P(Z)\to\mathcal P(X).
\]
\end{proposizione}
\begin{proof}
Per ogni $A\subseteq X$, $z\in g_*(f_*(A))$ significa che $z=g(y)$ per qualche $y=f(a)$ con $a\in A$: equivale a $z\in(g\circ f)_*(A)$.

Per ogni $C\subseteq Z$,
\[
x\in(g\circ f)^*(C)\iff g(f(x))\in C
\iff f(x)\in g^*(C)\iff x\in f^*(g^*(C)).
\]
La controimmagine inverte dunque l'ordine delle due applicazioni.
\end{proof}

\begin{proposizione}\label{prop:composte}
Per $f:X\to Y$ e $g:Y\to Z$:
\begin{enumerate}[label=\alph*)]
\item se $f,g$ sono iniettive, anche $g\circ f$ è iniettiva;
\item se $f,g$ sono suriettive, anche $g\circ f$ è suriettiva;
\item se $g\circ f$ è iniettiva, $f$ è iniettiva;
\item se $g\circ f$ è suriettiva, $g$ è suriettiva.
\end{enumerate}
\end{proposizione}
\begin{proof}
\textbf{a).} Da $g(f(x))=g(f(x'))$, l'iniettività di $g$ dà $f(x)=f(x')$; quella di $f$ dà $x=x'$.

\textbf{b).} Dato $z\in Z$, esiste $y\in Y$ con $g(y)=z$; esiste poi $x\in X$ con $f(x)=y$. Allora $(g\circ f)(x)=z$.

\textbf{c).} Se $f(x)=f(x')$, applicando $g$ si ottiene $(g\circ f)(x)=(g\circ f)(x')$. L'iniettività della composta dà $x=x'$.

\textbf{d).} Dato $z\in Z$, la suriettività della composta fornisce $x\in X$ con $g(f(x))=z$. Il valore $f(x)\in Y$ è una preimmagine di $z$ per $g$.
\end{proof}

\begin{esempio}[Quattro implicazioni che non valgono]
\textbf{$f$ iniettiva non basta per la composta.} Per $f,g:\R\to\R$, $f(x)=x^3$ e $g(y)=y^2$, la prima è iniettiva ma $(g\circ f)(x)=x^6$ non lo è, avendo lo stesso valore in $1$ e $-1$.

\textbf{La composta suriettiva non forza $f$ suriettiva.} Siano $f:\R\to\R$, $f(x)=x^2$, e $g:\R\to\R_{\ge0}$, $g(y)=|y|$. La composta è $x^2$ ed è suriettiva su $\R_{\ge0}$, ma $f$ non raggiunge i negativi del proprio codominio.

\textbf{La composta iniettiva non forza $g$ iniettiva.} L'inclusione $f:\R_{\ge0}\to\R$, seguita da $g:\R\to\R$, $g(y)=y^2$, dà una composta iniettiva. La funzione $g$ sull'intero dominio reale non lo è.

\textbf{$g$ suriettiva non basta per la composta.} L'inclusione $f:[0,1]\to\R$, seguita da $g:\R\to\R_{\ge0}$, $g(y)=y^2$, ha composta di immagine $[0,1]$, benché $g$ sia suriettiva.
\end{esempio}

\section{Inverse a sinistra e a destra}
% A–L L5 p. 7 e L6 pp. 1–5; M–Z L6 pp. 5–6 e L7 pp. 2–5.
L'\textbf{identità} di $X$ è $\operatorname{id}_X:X\to X$, $x\mapsto x$. Per ogni $f:X\to Y$, $\operatorname{id}_Y\circ f=f=f\circ\operatorname{id}_X$, perché i valori delle tre funzioni sono $f(x)$.

\begin{definizione}
Per $f:X\to Y$, una funzione $g:Y\to X$ è un'\textbf{inversa a sinistra} di $f$ se $g\circ f=\operatorname{id}_X$, ed è un'\textbf{inversa a destra} se $f\circ g=\operatorname{id}_Y$. La funzione $f$ è \textbf{invertibile} se ammette un'inversa da entrambi i lati, detta \textbf{inversa bilatera} e indicata con $f^{-1}$.
\end{definizione}

\begin{proposizione}[Unicità dell'inversa bilatera]
Se $g\circ f=\operatorname{id}_X$ e $f\circ h=\operatorname{id}_Y$, allora $g=h$. Quindi un'inversa a sinistra e una a destra, se esistono entrambe, coincidono e costituiscono un'inversa bilatera unica.
\end{proposizione}
\begin{proof}
Per associatività e proprietà dell'identità,
\[
g=g\circ\operatorname{id}_Y=g\circ(f\circ h)
=(g\circ f)\circ h=\operatorname{id}_X\circ h=h.
\]
La stessa uguaglianza prova l'unicità applicandola a due inverse bilatere.
\end{proof}

\begin{esempio}[Quadrato e radici]
Per $f:\R\to\R_{\ge0}$, $f(x)=x^2$, le funzioni $g_+(y)=\sqrt y$ e $g_-(y)=-\sqrt y$, da $\R_{\ge0}$ a $\R$, sono inverse a destra: $f(g_\pm(y))=y$. Sono distinte, per esempio in $y=1$. Invece $g_+(f(x))=|x|$ e $g_-(f(x))=-|x|$, perciò nessuna delle due è inversa a sinistra su tutto $\R$. Restringendo il dominio del quadrato a $\R_{\ge0}$ si ottiene l'inversa bilatera $\sqrt{\phantom{x}}$; restringendolo a $\R_{\le0}$ si ottiene $-\sqrt{\phantom{x}}$.
\end{esempio}

\begin{teorema}\label{teo:inverse}
Sia $f:X\to Y$ con $X\ne\varnothing$.
\begin{enumerate}[label=\alph*)]
\item $f$ è iniettiva se e solo se ammette un'inversa a sinistra.
\item $f$ è suriettiva se e solo se ammette un'inversa a destra, nel quadro con assioma della scelta.
\item $f$ è biettiva se e solo se è invertibile.
\end{enumerate}
\end{teorema}
\begin{proof}
\textbf{a), costruzione.} Supponiamo $f$ iniettiva. Per ogni $y\in\operatorname{Im}(f)$ esiste un'unica preimmagine $x_y$. Fissiamo $x_0\in X$ e definiamo su tutto $Y$
\[
g(y)=\begin{cases}x_y&y\in\operatorname{Im}(f),\\x_0&y\notin\operatorname{Im}(f).\end{cases}
\]
Il primo valore esiste ed è unico, il secondo è un elemento ammesso di $X$, dunque $g:Y\to X$ è ben definita. Per ogni $x\in X$, la preimmagine unica di $f(x)$ è $x$, quindi $g(f(x))=x$.

\textbf{a), necessità.} Se $g\circ f=\operatorname{id}_X$, la composta è iniettiva e la proposizione~\ref{prop:composte} implica che $f$ è iniettiva. Inoltre $g$ è suriettiva: ogni $x$ è $g(f(x))$.

\textbf{b), costruzione.} Supponiamo $f$ suriettiva. Ogni fibra $f^*(\{y\})$, $y\in Y$, è non vuota. Scegliendo simultaneamente un elemento $x_y$ in ciascuna fibra si pone $g(y)=x_y$. Questa è precisamente la scelta richiesta per una famiglia arbitraria di fibre. Per costruzione $f(g(y))=f(x_y)=y$, dunque $f\circ g=\operatorname{id}_Y$.

\textbf{b), necessità.} Se $f\circ g=\operatorname{id}_Y$, per ogni $y$ l'elemento $g(y)$ è una sua preimmagine, quindi $f$ è suriettiva. L'inversa destra $g$ è iniettiva: da $g(y)=g(y')$, applicando $f$, segue $y=y'$.

\textbf{c).} Se $f$ è biettiva, ogni $y$ ha un'unica preimmagine $x_y$: ponendo $g(y)=x_y$ si verificano entrambe le identità precedenti, senza scelte arbitrarie. Viceversa un'inversa bilatera è sia inversa a sinistra sia inversa a destra, e i due versi necessari già provati danno iniettività e suriettività.
\end{proof}

\begin{osservazione}[Il dominio vuoto]
L'ipotesi $X\ne\varnothing$ serve in a) per il valore fuori dall'immagine. La funzione vuota $\varnothing\to Y$ è iniettiva, ma per $Y\ne\varnothing$ non esiste alcuna funzione $Y\to\varnothing$. Per $X=Y=\varnothing$ la funzione vuota è biettiva e inversa di se stessa. La definizione generale ammette questi casi, anche se le costruzioni precedenti sono state formulate per domini non vuoti.
\end{osservazione}

Per una biezione, $f^{-1}(y)=x$ equivale a $f(x)=y$: il grafico dell'inversa si ottiene scambiando le due coordinate del grafico di $f$. Nel piano è la riflessione rispetto alla retta $y=x$. La funzione del piano studiata sopra ha quindi inversa
\[
F^{-1}(u,v)=\left(\frac{v}{1+(\sqrt[3]u)^2},\sqrt[3]u\right).
\]

\begin{corollario}
Per insiemi non vuoti $X,Y$, esiste un'iniezione $X\to Y$ se e solo se esiste una suriezione $Y\to X$, nel quadro con assioma della scelta.
\end{corollario}
\begin{proof}
Un'iniezione $X\to Y$ ha un'inversa a sinistra $Y\to X$, che è suriettiva. Viceversa una suriezione $Y\to X$ ha un'inversa a destra $X\to Y$, che è iniettiva. Sono esattamente i due versi richiesti.
\end{proof}

Si scrive $X\hookrightarrow Y$ per un'iniezione, $X\twoheadrightarrow Y$ per una suriezione e $X\xrightarrow{\sim}Y$ per una biezione. Se $X\subseteq Y$, l'\textbf{inclusione} è $i:X\hookrightarrow Y$, $i(x)=x$.

\begin{esercizio}[Inverse destre del seno]
Trovare due inverse a destra distinte di $s:\R\to[-1,1]$, $s(x)=\sin x$. Può esistere un'inversa a sinistra?
\end{esercizio}
\begin{svolgimento}
Si assume la proprietà analitica usuale che la restrizione del seno a $[-\pi/2,\pi/2]$ sia una biezione su $[-1,1]$, con inversa $\arcsin$. Ponendo $g_1(y)=\arcsin y$ e $g_2(y)=\pi-\arcsin y$, si ottiene
\[
\sin(g_1(y))=y,\qquad \sin(g_2(y))=\sin(\arcsin y)=y.
\]
Le due funzioni sono distinte perché $g_1(0)=0$ e $g_2(0)=\pi$. Non esiste un'inversa a sinistra, poiché $s(0)=s(\pi)$ rende $s$ non iniettiva.
\end{svolgimento}

\section{Insiemi di funzioni e caratteristiche}
% A–L L6 pp. 6–7 e L7 p. 1; M–Z L7 pp. 5–9.
\begin{definizione}
L'insieme di tutte le funzioni da $X$ a $\Omega$ si indica con $\Omega^X$. Per $X=\{1,\ldots,n\}$, una funzione è determinata dalla $n$-upla $(\varphi(1),\ldots,\varphi(n))$, quindi $\Omega^X$ si identifica con $\Omega^n$. Per esempio $\R^{\{1,2,3,4\}}$ si identifica con $\R^4$.
\end{definizione}

Una successione reale $(x_0,x_1,\ldots)$ è una funzione $\N\to\R$, $i\mapsto x_i$, dunque appartiene a $\R^\N$. Più in generale una famiglia $(x_i)_{i\in I}$ di elementi di $\Omega$ è una funzione $I\to\Omega$: a ogni indice è assegnato un unico valore, mentre indici distinti possono avere lo stesso valore.

\begin{definizione}
Per $A\subseteq X$, la \textbf{funzione caratteristica} è
\[
\chi_A:X\to\{0,1\},\qquad \chi_A(x)=\begin{cases}1&x\in A,\\0&x\notin A.\end{cases}
\]
\end{definizione}

\begin{teorema}\label{teo:caratteristiche}
Le applicazioni
\[
\begin{aligned}
\alpha:\{0,1\}^X&\to\mathcal P(X),&\alpha(\varphi)&=\varphi^*(\{1\}),\\
\beta:\mathcal P(X)&\to\{0,1\}^X,&\beta(A)&=\chi_A
\end{aligned}
\]
sono biezioni inverse.
\end{teorema}
\begin{proof}
Le applicazioni sono ben definite: $\varphi^*(\{1\})$ è un sottoinsieme di $X$ e $\chi_A$ è una funzione a valori in $\{0,1\}$.

\passomath{\alpha\circ\beta=\operatorname{id}_{\mathcal P(X)}} Per ogni $A\subseteq X$,
\[
\alpha(\beta(A))=\{x\in X\mid\chi_A(x)=1\}=A.
\]
\passomath{\beta\circ\alpha=\operatorname{id}_{\{0,1\}^X}} Fissata $\varphi:X\to\{0,1\}$, la caratteristica di $\varphi^*(\{1\})$ vale $1$ dove $\varphi(x)=1$ e vale $0$ negli altri punti. In questi ultimi $\varphi(x)=0$, perché non sono ammessi altri valori. Le due funzioni coincidono in ogni punto.

Le identità danno inverse bilatere e quindi biezioni. Per $X=\varnothing$ l'argomento resta valido: sia $\mathcal P(X)$ sia $\{0,1\}^X$ hanno un unico elemento.
\end{proof}

% !TeX root = ../main.tex
\chapter{Insiemi finiti, infiniti e cardinalità}
\fonti{Lezioni 7, 8 e 9}

\section{Insiemi finiti e infiniti}
% A–L L7 p. 2; M–Z L8 pp. 1–2. Le proprietà elementari finite sono assunte.
\begin{definizione}
Poniamo $[n]=\{1,\ldots,n\}$ per $n\ge1$ e $[0]=\varnothing$. Un insieme $X$ è \textbf{finito} se esiste una biezione $[n]\to X$ per qualche $n\in\N$; è \textbf{infinito} se non è finito. Nel caso finito il numero $n$, univocamente determinato, si chiama \textbf{cardinalità} e si indica con $|X|$.
\end{definizione}

Assumiamo i fatti elementari sul conteggio finito: $[m]$ e $[n]$ sono in biezione solo se $m=n$; un sottoinsieme di un insieme finito è finito; un'iniezione di un insieme finito in se stesso è suriettiva. In particolare un insieme finito non può essere in biezione con un proprio sottoinsieme stretto.

\begin{corollario}
Se $X$ è finito, $|\mathcal P(X)|=2^{|X|}$.
\end{corollario}
\begin{proof}
Per il teorema~\ref{teo:caratteristiche}, i sottoinsiemi di $X$ corrispondono alle funzioni $X\to\{0,1\}$. Se $|X|=n$, per ciascuno degli $n$ elementi si sceglie indipendentemente uno dei due valori: le scelte sono $2^n$. Per $n=0$ c'è la sola funzione vuota e $\mathcal P(\varnothing)=\{\varnothing\}$ ha $1=2^0$ elemento.
\end{proof}

\begin{teorema}\label{teo:infiniti}
Nel quadro con assioma della scelta, per un insieme $X$ sono equivalenti:
\begin{enumerate}[label=\alph*)]
\item $X$ è infinito;
\item esiste un'iniezione $f:\N\to X$;
\item esiste un'iniezione $g:X\to X$ non suriettiva;
\item esiste una biezione fra $X$ e un suo sottoinsieme proprio.
\end{enumerate}
\end{teorema}
\begin{proof}
% A–L L7 pp. 2–3 e L8 pp. 1–2, più dettagliata sui tre casi dello spostamento.
\passomath{a\Rightarrow b} Poiché $X$ è infinito, è non vuoto: scegliamo $x_0\in X$. Scelti $x_0,\ldots,x_{n-1}$ distinti, il complementare $X\setminus\{x_0,\ldots,x_{n-1}\}$ è non vuoto, altrimenti $X$ sarebbe finito. Possiamo dunque scegliere $x_n$ in esso. L'assioma della scelta garantisce la successione di scelte per tutti gli $n$. La funzione $f(n)=x_n$ è iniettiva: per $m<n$, $x_n$ è stato scelto diverso da $x_m$.

\passomath{b\Rightarrow c} Posto $S=f(\N)$, definiamo
\[
g(x)=\begin{cases}f(n+1)&x=f(n)\in S,\\x&x\notin S.\end{cases}
\]
L'indice $n$ è unico perché $f$ è iniettiva; la regola è quindi ben definita. Per l'iniettività distinguiamo i casi. Se $x=f(n)$ e $x'=f(m)$, da $g(x)=g(x')$ segue $f(n+1)=f(m+1)$, dunque $n=m$ e $x=x'$. Se entrambi sono fuori da $S$, sono lasciati fissi, quindi immagini uguali impongono punti uguali. Se solo uno è in $S$, anche solo una delle immagini è in $S$: non possono coincidere. Infine $f(0)$ non è raggiunto: gli elementi fuori da $S$ rimangono fuori, e quelli in $S$ vengono mandati nei punti $f(n+1)\ne f(0)$. Dunque $g$ non è suriettiva.

\passomath{c\Rightarrow d} Restringendo il codominio di $g$ a $g(X)$ si ottiene una biezione $X\to g(X)$. La mancata suriettività dice esattamente $g(X)\subsetneq X$.

\passomath{d\Rightarrow a} Una biezione $h:X\to S\subsetneq X$, composta con l'inclusione $S\hookrightarrow X$, dà un'iniezione non suriettiva $X\to X$. Questo è impossibile per un insieme finito, per il fatto di conteggio assunto sopra. Quindi $X$ è infinito.
\end{proof}

\begin{schema}
\begin{tikzpicture}[>=Stealth]
  \foreach \i in {0,...,3}{\node (x\i) at (2*\i,1) {$f(\i)$};\node (y\i) at (2*\i,-.5) {$f(\i)$};}
  \foreach \i in {0,...,2}{\pgfmathtruncatemacro{\j}{\i+1}\draw[->,coloreschema] (x\i)--(y\j);}
  \node at (8,1) {$\cdots$};\node at (8,-.5) {$\cdots$};\draw[->,coloreschema] (x3)--(7.7,-.3);
  \node[below=2pt] at (y0.south) {omesso};\node at (4,-1.6) {fuori da $f(\N)$ ogni elemento resta fisso};
\end{tikzpicture}
\end{schema}

\section{Equipotenza e confronto di cardinalità}
% A–L L8 pp. 2–3; M–Z L8 pp. 5–6.
\begin{definizione}
Due insiemi sono \textbf{equipotenti}, e si scrive $X\simeq Y$ oppure $|X|=|Y|$, se esiste una biezione $X\to Y$. Si scrive $|X|\le|Y|$ se esiste un'iniezione $X\to Y$; si scrive $|X|<|Y|$ se vale $\le$ ma non l'uguaglianza.
\end{definizione}

L'equipotenza è riflessiva, tramite l'identità; simmetrica, tramite l'inversa di una biezione; transitiva, tramite la composizione di biezioni. La notazione $|X|$ per insiemi arbitrari esprime la grandezza comune agli insiemi equipotenti, senza richiedere qui una costruzione formale dei cardinali. Per esempio $|\{\text{rosso},\text{blu},\text{verde}\}|=3$.

Per insiemi non vuoti, $|X|\le|Y|$ equivale all'esistenza di una suriezione $Y\to X$, per il corollario sui criteri delle inverse. La formulazione tramite iniezioni vale anche per il vuoto: $|\varnothing|\le|Y|$ sempre, benché non esista una funzione $Y\to\varnothing$ se $Y\ne\varnothing$.

\begin{teorema}[Cantor--Schröder--Bernstein]\label{teo:csb}
Se esistono iniezioni $X\to Y$ e $Y\to X$, allora esiste una biezione $X\to Y$. Dunque $|X|\le|Y|$ e $|Y|\le|X|$ implicano $|X|=|Y|$.
\end{teorema}
Il teorema è assunto senza dimostrazione. Il confronto delle cardinalità è riflessivo e transitivo perché l'identità è iniettiva e la composizione di iniezioni è iniettiva; l'antisimmetria è precisamente il risultato appena enunciato.

\begin{teorema}[Confronto totale]
Nel quadro con assioma della scelta, per ogni coppia di insiemi $X,Y$ vale $|X|\le|Y|$ oppure $|Y|\le|X|$.
\end{teorema}
Anche questo risultato è assunto senza dimostrazione. In particolare, per il teorema~\ref{teo:infiniti}, la cardinalità di $\N$ è la più piccola cardinalità infinita.

\section{Numerabilità}
\begin{definizione}
Un insieme $X$ è \textbf{numerabile} se $|X|\le|\N|$; è \textbf{numerabile infinito} se $|X|=|\N|$. La convenzione include fra i numerabili gli insiemi finiti e il vuoto. Scriviamo $\aleph_0=|\N|$.
\end{definizione}

Un insieme numerabile infinito è in biezione con $\N$: se $X\hookrightarrow\N$ e $X$ è infinito, il teorema~\ref{teo:infiniti} dà anche $\N\hookrightarrow X$, e si applica Cantor--Schröder--Bernstein.

\begin{proposizione}
$|\N\times\N|=\aleph_0$ e $|\N^r|=\aleph_0$ per ogni intero $r\ge1$.
\end{proposizione}
\begin{proof}
\textbf{Coppie.} Consideriamo $p:\N\times\N\to\N$, $p(m,n)=2^m(2n+1)-1$. Se $p(m,n)=p(m',n')$, allora
\[
2^m(2n+1)=2^{m'}(2n'+1).
\]
Se $m<m'$, cancellando $2^m$ si avrebbe un numero dispari uguale a uno pari, impossibile. Lo stesso esclude $m>m'$ scambiando i ruoli. Quindi $m=m'$, e cancellando la potenza di $2$ si ricava $n=n'$. L'applicazione è iniettiva; nell'altro verso $m\mapsto(m,0)$ è un'iniezione. Cantor--Schröder--Bernstein dà l'equipotenza. Non si è usata la fattorizzazione unica in primi.

\textbf{Prodotti finiti, induzione su $r$.} Per $r=1$ vale l'identità. Supponiamo data una biezione $u_r:\N^r\to\N$. L'applicazione $(a_1,\ldots,a_{r+1})\mapsto(u_r(a_1,\ldots,a_r),a_{r+1})$ è una biezione $\N^{r+1}\to\N^2$: l'inversa applica $u_r^{-1}$ alla prima coordinata. Componendola con una biezione $\N^2\to\N$ si ottiene il passo $r+1$. Il principio d'induzione conclude.
\end{proof}

\begin{osservazione}[Enumerazione per diagonali]
Un secondo metodo consiste nel percorrere successivamente le diagonali $m+n=0,1,2,\ldots$. Ciascuna ha un numero finito di punti e ogni coppia appartiene a una sola diagonale. Elencando ogni diagonale una volta si ottiene una successione senza ripetizioni che raggiunge tutte le coppie. Questo dà direttamente una biezione $\N\to\N^2$.
\end{osservazione}

\begin{schema}
\begin{tikzpicture}[scale=.55,>=Stealth]
  \draw[asse] (-.4,0)--(5.5,0) node[right] {$m$};\draw[asse] (0,-.4)--(0,5.5) node[above] {$n$};
  \foreach \x in {0,...,4}{\foreach \y in {0,...,4}{\fill (\x,\y) circle (1.7pt);}}
  \draw[->,coloreschema,thick] (0,0)--(1,0)--(0,1)--(0,2)--(1,1)--(2,0)--(3,0)--(2,1)--(1,2)--(0,3);
  \node[right,align=left] at (6,2.5) {Un'altra enumerazione percorre le diagonali\\$m+n=0,1,2,\ldots$ del reticolo.};
\end{tikzpicture}
\end{schema}

\begin{teorema}[Unione numerabile di numerabili]\label{teo:unione-numerabile}
Se $I$ è numerabile e ogni $A_i$, $i\in I$, è numerabile, allora $\bigcup_{i\in I}A_i$ è numerabile, nel quadro con assioma della scelta.
\end{teorema}
\begin{proof}
\textbf{Riduzione ai membri non vuoti.} Se l'unione è vuota, la conclusione vale. Altrimenti poniamo $I'=\{i\in I\mid A_i\ne\varnothing\}$. È numerabile e non vuoto, quindi esiste una suriezione $e:\N\to I'$: se $I'$ è infinito si usa una biezione, se è finito si ripete un elenco dei suoi elementi.

\textbf{Costruzione della suriezione.} Ogni $A_i$ non vuoto è a sua volta immagine di una suriezione $\varphi_i:\N\to A_i$. Scegliamo una tale funzione per ciascun $i\in I'$: è qui che occorre la scelta simultanea. Definiamo
\[
\Phi:\N^2\to\bigcup_{i\in I}A_i,\qquad \Phi(k,n)=\varphi_{e(k)}(n).
\]
Per ogni $a$ dell'unione, esiste $i\in I'$ con $a\in A_i$; esiste $k$ con $e(k)=i$ e poi $n$ con $\varphi_i(n)=a$. Quindi $\Phi$ è suriettiva. Componendo con una biezione $\N\to\N^2$ si ha una suriezione da $\N$ sull'unione, dunque un'iniezione dell'unione in $\N$ per il criterio già provato. L'unione è numerabile.
\end{proof}

\begin{corollario}
$|\Z|=|\Q|=\aleph_0$.
\end{corollario}
\begin{proof}
\textbf{Interi.} Le formule $e(2k)=k$ ed $e(2k+1)=-(k+1)$ definiscono una biezione $\N\to\Z$: i due rami raggiungono rispettivamente tutti i non negativi e tutti i negativi, senza ripetizioni.

\textbf{Razionali.} Il prodotto $\Z\times\N^*$ è numerabile perché si inietta in $\N^2$, usando una biezione $\Z\to\N$ e l'inclusione $\N^*\subseteq\N$. La funzione $(a,b)\mapsto a/b$ è suriettiva su $\Q$: ogni razionale ammette un denominatore positivo cambiando entrambi i segni, se necessario. Dunque $|\Q|\le\aleph_0$. L'inclusione $\N\hookrightarrow\Q$ dà la disuguaglianza opposta e si conclude per Cantor--Schröder--Bernstein.
\end{proof}

\section{Cantor e la cardinalità del continuo}
% A–L L8 pp. 5–6; M–Z L8 pp. 2–3. Nessun anticipo di Cantor prima degli infiniti.
\begin{teorema}[Cantor]\label{teo:cantor}
Per ogni insieme $X$ non esiste una suriezione $X\to\mathcal P(X)$. In particolare $|X|<|\mathcal P(X)|$.
\end{teorema}
\begin{proof}
\textbf{Diagonale.} Data una qualunque $f:X\to\mathcal P(X)$, definiamo $D=\{x\in X\mid x\notin f(x)\}$. È un sottoinsieme di $X$, dunque un elemento del codominio. Se $D=f(a)$, allora
\[
a\in D\iff a\notin f(a)\iff a\notin D.
\]
Se $a\in D$ si ottiene $a\notin D$; se $a\notin D$ si ottiene $a\in D$. Entrambi i casi sono impossibili. Quindi $D$ non è nell'immagine e $f$ non è suriettiva.

\textbf{Confronto stretto.} La funzione $x\mapsto\{x\}$ è iniettiva, quindi $|X|\le|\mathcal P(X)|$. Non può esserci una biezione, perché sarebbe una suriezione. Per $X=\varnothing$, la conclusione vale direttamente: $\mathcal P(X)=\{\varnothing\}$ ha un elemento.
\end{proof}

\begin{definizione}
Si pone $2^{|X|}=|\{0,1\}^X|=|\mathcal P(X)|$. La \textbf{cardinalità del continuo} è $\mathfrak c=2^{\aleph_0}$. Il teorema di Cantor dà $\aleph_0<\mathfrak c$ e, iterato, dà
\[
\aleph_0<2^{\aleph_0}<2^{\,2^{\aleph_0}}<\cdots.
\]
\end{definizione}

\begin{osservazione}[Ipotesi del continuo]
L'\textbf{ipotesi del continuo} afferma che non esiste una cardinalità strettamente compresa fra $\aleph_0$ e $2^{\aleph_0}$. Non è una conseguenza del solo teorema di Cantor. Si assume qui, senza dimostrarlo, il risultato di logica secondo cui essa è indipendente dagli assiomi usuali degli insiemi con scelta, se questi sono coerenti. L'ipotesi stessa non viene assunta per i risultati di questa dispensa.
\end{osservazione}

\begin{esercizio}
Siano $A$ finito e $B$ numerabile. Sono numerabili $A\cup B$, $A\times B$ e $B^A$?
\end{esercizio}
\begin{svolgimento}
\textbf{Unione.} Un insieme finito è numerabile. Il teorema~\ref{teo:unione-numerabile}, applicato alla famiglia di due membri $A,B$, dà la numerabilità dell'unione.

\textbf{Prodotto.} Iniezioni $A\to\N$ e $B\to\N$ inducono, coordinata per coordinata, un'iniezione $A\times B\to\N^2$. Quindi il prodotto è numerabile.

\textbf{Funzioni.} Per $A=\{a_1,\ldots,a_m\}$, $m\ge1$, la mappa $f\mapsto(f(a_1),\ldots,f(a_m))$ è una biezione $B^A\to B^m$, perché ogni scelta della $m$-upla assegna esattamente una funzione. L'iniezione $B\to\N$ dà $B^m\hookrightarrow\N^m$, che è numerabile. Se $A=\varnothing$, $B^A$ ha la sola funzione vuota; se $B=\varnothing$ e $A\ne\varnothing$, non esistono funzioni e $B^A=\varnothing$. Anche nei casi vuoti tutte e tre le risposte sono affermative.
\end{svolgimento}

\section{Cardinalità della retta reale}
% A–L L9 p. 1; M–Z L9 pp. 2–3. Passo analitico assunto, come dichiarato in M–Z.
\begin{proposizione}
Una biezione $u:X\to Y$ induce una biezione $u_*:\mathcal P(X)\to\mathcal P(Y)$.
\end{proposizione}
\begin{proof}
La candidata inversa è $(u^{-1})_*$. Le regole di composizione delle immagini dirette danno $(u^{-1})_*\circ u_*=(u^{-1}\circ u)_*=(\operatorname{id}_X)_*$, che è l'identità di $\mathcal P(X)$. Nell'altro verso, $u_*\circ(u^{-1})_*=(\operatorname{id}_Y)_*$, identità di $\mathcal P(Y)$. Dunque le due applicazioni sono inverse.
\end{proof}

\begin{teorema}\label{teo:continuo}
$|\R|=\mathfrak c=2^{\aleph_0}$.
\end{teorema}
\begin{proof}[Dimostrazione con un risultato analitico assunto]
\passomath{|\R|\le\mathfrak c} A ogni $x\in\R$ associamo il taglio razionale inferiore $L_x=\{q\in\Q\mid q<x\}$. Se $x<y$, per densità esiste $q\in\Q$ con $x<q<y$: esso appartiene a $L_y$, ma non a $L_x$. La mappa $x\mapsto L_x$ è quindi iniettiva. Poiché $\Q\simeq\N$, la proposizione precedente dà $|\mathcal P(\Q)|=|\mathcal P(\N)|=\mathfrak c$, dunque $|\R|\le\mathfrak c$.

\passomath{\mathfrak c\le|\R|} Si assume il risultato analitico che esiste un'iniezione $\{0,1\}^{\N}\to\R$.\footnote{La costruzione mediante serie numeriche richiede un argomento aggiuntivo non sviluppato qui. La formula binaria da sola non prova l'iniettività: $(1,0,0,\ldots)$ e $(0,1,1,\ldots)$ danno entrambe $1$ nella somma $\sum_{n\ge0}\varepsilon(n)/2^n$.} Da questo fatto e dalla prima disuguaglianza segue l'uguaglianza per Cantor--Schröder--Bernstein.
\end{proof}

\section{Somme e prodotti di cardinalità infinite}
% Integrazione M–Z L9 pp. 4–5, espressamente autorizzata dall'autore.
\begin{definizione}
L'\textbf{unione disgiunta} è $A\sqcup B=(A\times\{0\})\cup(B\times\{1\})$: le etichette mantengono distinte le copie anche se $A\cap B\ne\varnothing$. Per $|A|=\alpha$, $|B|=\beta$, si definiscono $\alpha+\beta=|A\sqcup B|$ e $\alpha\beta=|A\times B|$.
\end{definizione}

Le definizioni non dipendono dagli insiemi che rappresentano le cardinalità. Infatti biezioni $u:A\to A'$ e $v:B\to B'$ inducono, sul prodotto, $(a,b)\mapsto(u(a),v(b))$, con inversa coordinata per coordinata; sull'unione disgiunta si usano $(a,0)\mapsto(u(a),0)$ e $(b,1)\mapsto(v(b),1)$, ancora invertibili.

\begin{teorema}
Nel quadro con assioma della scelta, se $\alpha,\beta$ sono cardinalità non nulle e almeno una è infinita,
\[
\alpha+\beta=\alpha\beta=\max\{\alpha,\beta\}.
\]
\end{teorema}
Il risultato generale è assunto senza dimostrazione. L'ipotesi di fattori non vuoti è necessaria per il prodotto: se $B=\varnothing$, $A\sqcup B\simeq A$, ma $A\times B=\varnothing$.

% !TeX root = ../main.tex
\chapter{Aritmetica degli interi}\label{cap:interi}
\fonti{Lezioni 9, 10 e 11}

\section{Interi e unità}
% A–L L9 pp. 2–3; M–Z L9 pp. 5–6.
Assumiamo le proprietà usuali delle operazioni interne $+,\cdot:\Z^2\to\Z$. L'addizione è associativa e commutativa, ha neutro $0$ e ogni $z$ ha opposto $-z$; la moltiplicazione è associativa e commutativa, ha neutro $1$ ed è distributiva rispetto all'addizione: $(a+b)c=ac+bc$. Queste proprietà fanno di $\Z$ un \textbf{anello commutativo con unità}.

Inoltre non vi sono divisori dello zero: $ab=0$ implica $a=0$ oppure $b=0$. Ne segue la cancellazione: se $a\ne0$ e $ab=ac$, allora $a(b-c)=0$ e perciò $b=c$.

\begin{definizione}
Un intero $a$ è un'\textbf{unità}, o è \textbf{invertibile}, se esiste $b\in\Z$ con $ab=1$.
\end{definizione}

Si ha $U(\Z)=\{1,-1\}$.: entrambi sono inversi di se stessi; viceversa $ab=1$ implica $|a||b|=1$, quindi i due valori assoluti interi positivi sono entrambi $1$ \footnote{La dimostrazione non é del tutto triviale: posso ragionare sui moduli ma prima devo capire come si comporta la moltiplicazione rispetto al $\leq$.}

\section{Divisibilità}
\begin{definizione}
Per $a,b\in\Z$ si dice che $a$ \textbf{divide} $b$, e si scrive $a\mid b$, se $b=am$ per qualche $m\in\Z$. Inoltre, se $a\mid b$ e $b\mid a$, si dice che $a$ e $b$ sono \textbf{associati} e si scrive $a\sim b$.
\end{definizione}

\begin{proposizione}
La divisibilità è riflessiva e transitiva, ma non è né simmetrica né antisimmetrica (quindi in $\Z$ non è una relazione di ordine o di equivalenza\footnote{E in $\N$? Beh è riflessiva, antisym. e transitiva (ma non totale) quindi è un ordine parziale.}). Due interi sono associati se e solo se $a=\pm b$; $\sim$ è un'equivalenza.
\end{proposizione}
\begin{proof}
Si ha $a=a\cdot1$; se $b=am$ e $c=bn$, allora $c=a(mn)$. La simmetria fallisce perché $2\mid4$ e $4\nmid2$; l'antisimmetria fallisce perché $1\mid-1$ e $-1\mid1$.

Se $b=am$ e $a=bn$, per $a=0$ segue $b=0$. Se $a\ne0$, sostituendo si ha $a=amn$; cancellando $a$ si ricava $mn=1$, quindi $m=n=1$ oppure $m=n=-1$. Dunque $b=\pm a$. Il viceversa segue moltiplicando per $1$ o $-1$. Infine l'uguaglianza a meno del segno è riflessiva, simmetrica e transitiva: se $a=\varepsilon b$ e $b=\eta c$, con $\varepsilon,\eta\in\{1,-1\}$, allora $a=(\varepsilon\eta)c$.
\end{proof}

\begin{osservazione}
Per ogni $z\in\Z$, $1,-1,z,-z$ dividono $z$. Ogni intero divide $0$, perché $0=z\cdot0$; invece $0$ divide soltanto $0$. La divisibilità non è totale: $2\nmid3$ e $3\nmid2$.
\end{osservazione}

\section{Primi e crivello di Eratostene}
% A–L L9 pp. 3–4. L'infinità precede MCD e divisione, senza usare fattorizzazione unica.
\begin{definizione}
Un intero $p$ è \textbf{primo} se $p\ne0,\pm1$ e da $p=ab$ segue $a=\pm1$ oppure $b=\pm1$. Equivalentemente, i divisori di $p$ sono soltanto $\pm1,\pm p$. Indichiamo con $P$ l'insieme dei primi positivi; $p$ è primo se e solo se lo è $-p$.\footnote{La "vera" definizione di primo è "$p$ è primo $\iff p\ne0,\pm1$ e $p\mid ab\Rightarrow p\mid a$ oppure $p\mid b$" (questa infatti si generalizza). Dopo si dimostra che sono equivalenti.}
\end{definizione}

\begin{teorema}
Valgono le seguenti:
\begin{enumerate}
    \item Ogni intero $N>1$ ha un divisore primo;
    \item I primi positivi sono infiniti e, in particolare, $|P|=\aleph_0$.
\end{enumerate}
\end{teorema}
\begin{proof}
\textbf{1.} Ogni intero $N>1$ ha un divisore primo: l'insieme dei suoi divisori positivi maggiori di $1$ è non vuoto, perché contiene $N$. Sia $q$ il minimo. Se $q$ non fosse primo, avrebbe un divisore positivo $u$ con $1<u<q$. Per transitività $u\mid N$, contro la minimalità di $q$.

\textbf{2.} I primi positivi sono infiniti: supponiamo per assurdo che siano finiti, e li ordiniamo in una lista $P=\{p_1,\ldots,p_r\}$. Consideriamo $N=p_1\cdots p_r+1>1$ e scegliamo un suo divisore primo $q$. Per ipotesi $q=p_i$ per qualche $i$, quindi divide anche il prodotto. Scrivendo $N=qu$ e $p_1\cdots p_r=qv$, si ha $1=q(u-v)$: un primo dividerebbe $1$, impossibile. Perciò non esiste una lista finita di tutti i primi.

Essendo un sottoinsieme infinito di $\N$, $P$ è numerabile infinito.
\end{proof}

\begin{osservazione}
Il numero $p_1\cdots p_r+1$ non deve essere primo: deve soltanto avere un divisore primo, che non può stare nella lista. La prova non richiede il teorema di fattorizzazione unica.
\end{osservazione}

\begin{esempio}[Crivello do Eratostene]
Per trovare i primi fino a un intero $N\ge2$, si elencano i numeri da $2$ a $N$, si conserva il più piccolo non cancellato e si cancellano i suoi multipli maggiori di esso, ripetendo il procedimento. Un primo non viene cancellato; un composto viene cancellato quando si esamina un suo divisore primo. Per $N=20$ rimangono $2,3,5,7,11,13,17,19$.

Si può interrompere la cancellazione dopo i primi non maggiori di $\sqrt N$: se $n\le N$ è composto, $n=uv$ con $u,v>1$ e almeno uno dei due fattori non supera $\sqrt n\le\sqrt N$. Quel fattore ha a sua volta un divisore primo non maggiore di $\sqrt N$, che fa cancellare $n$.
\end{esempio}

\begin{schema}
\begin{tikzpicture}[x=.76cm,y=.68cm]
\foreach \n in {1,...,20}{\pgfmathtruncatemacro{\xx}{mod(\n-1,10)}\pgfmathtruncatemacro{\yy}{-floor((\n-1)/10)}\node at (\xx,\yy) {$\n$};}
\foreach \n in {2,3,5,7,11,13,17,19}{\pgfmathtruncatemacro{\xx}{mod(\n-1,10)}\pgfmathtruncatemacro{\yy}{-floor((\n-1)/10)}\draw[coloreschema] (\xx,\yy) circle (.28);}
\foreach \n in {1,4,6,8,9,10,12,14,15,16,18,20}{\pgfmathtruncatemacro{\xx}{mod(\n-1,10)}\pgfmathtruncatemacro{\yy}{-floor((\n-1)/10)}\draw[black!50] (\xx-.19,\yy-.20)--(\xx+.19,\yy+.20);}
\end{tikzpicture}
\end{schema}

\section{Massimo comune divisore e minimo comune multiplo}
% A–L L9 pp. 4–6 e L10 p. 1; M–Z L9 p. 7. Prima definizioni, poi esistenza.
\begin{definizione}
Per $(a,b)\ne(0,0)$, un intero $d$ è un \textbf{massimo comune divisore} nel senso della divisibilità se $d\mid a$, $d\mid b$ e ogni divisore comune $c$ di $a,b$ divide $d$. Quando esiste, quello positivo si indica con $\operatorname{MCD}(a,b)$.
\end{definizione}

\begin{esempio}
Gli elenchi dei divisori sono
\[
\begin{aligned}
-12&:\quad \pm1,\pm2,\pm3,\pm4,\pm6,\pm12,\\
16&:\quad \pm1,\pm2,\pm4,\pm8,\pm16.
\end{aligned}
\]
I divisori comuni sono $\pm1,\pm2,\pm4$. Il numero $4$ divide entrambi e viene diviso da tutti i divisori comuni, quindi $\operatorname{MCD}(-12,16)=4$.
\end{esempio}

\begin{proposizione}
Se $d$ è un massimo comune divisore, anche $-d$ lo è. Due massimi comuni divisori sono associati: ne esiste al più uno positivo. Per $a\ne0$, $\operatorname{MCD}(a,0)=\operatorname{MCD}(0,a)=|a|$.
\end{proposizione}

\begin{proof}
Il segno non modifica i divisori né i multipli di un numero, quindi le due condizioni della definizione valgono per $d$ se e solo se valgono per $-d$. Se $d,d'$ soddisfano entrambe la definizione, ciascuno divide l'altro per la condizione universale: sono associati, dunque uguali a meno del segno. Non sono nulli, perché almeno uno tra $a,b$ è non nullo e $0$ non lo divide. Pertanto i rappresentanti positivi coincidono.

Il numero $|a|$ divide $a$ e $0$; ogni divisore comune divide $a$ e quindi $|a|$. Soddisfa così la definizione per $(a,0)$; lo scambio degli argomenti non cambia le condizioni.
\end{proof}

L'esistenza del massimo comune divisore per ogni coppia non nulla sarà ottenuta dall'algoritmo di Euclide. Fino a quel punto i risultati che lo menzionano si intendono per un massimo comune divisore esistente.

\begin{definizione}
Un intero $m$ è un \textbf{minimo comune multiplo} di $a,b$ nel senso della divisibilità se $a\mid m$, $b\mid m$ e $m$ divide ogni multiplo comune di $a,b$. Per $ab\ne0$, quello positivo, quando esiste, si indica con $\operatorname{mcm}(a,b)$. Se un argomento è zero, si pone $\operatorname{mcm}(a,b)=0$.
\end{definizione}

\begin{proposizione}
Cambiare il segno di un minimo comune multiplo dà ancora un minimo comune multiplo; due di essi sono associati. Se $ab\ne0$, ne esiste al più uno positivo. Se un argomento è zero, l'unico multiplo comune è $0$, che soddisfa la definizione.
\end{proposizione}
\begin{proof}
Il cambio di segno conserva le condizioni di divisibilità. Se $m,m'$ soddisfano la definizione, la condizione universale dà sia $m\mid m'$ sia $m'\mid m$. Quando $ab\ne0$, un minimo comune multiplo non può essere zero: il multiplo comune $ab\ne0$ dovrebbe essere divisibile per $0$. Quindi due rappresentanti positivi coincidono.

Se $a=0$, la condizione $0\mid m$ forza $m=0$; il numero $0$ è multiplo di $b$ e divide l'unico multiplo comune, che è ancora $0$. Lo stesso vale se $b=0$.
\end{proof}

\begin{osservazione}
«Massimo» e «minimo» si riferiscono alla divisibilità, non inizialmente all'ordine usuale. Per il minimo comune multiplo positivo è necessario $ab\ne0$: la sola esclusione $(a,b)=(0,0)$ non basta, come mostra la coppia $(0,2)$.
\end{osservazione}

\section{Combinazioni intere e coprimalità}
% A–L L10 p. 1; M–Z L9 pp. 8–9.
\begin{proposizione}\label{prop:combinazioni}
Se $d\mid a$ e $d\mid b$, allora $d\mid\alpha a+\beta b$ per ogni $\alpha,\beta\in\Z$. Se invece $d\mid a$ e $d'\mid b$, allora $dd'\mid ab$.
\end{proposizione}
\begin{proof}
Nella prima affermazione, $a=du$ e $b=dv$ danno $\alpha a+\beta b=d(\alpha u+\beta v)$. Nella seconda, $a=du$ e $b=d'v$ danno $ab=dd'(uv)$. I quozienti esibiti sono interi.
\end{proof}

\begin{proposizione}[Criterio tramite una combinazione]\label{prop:criterio-mcd}
Siano $(a,b)\ne(0,0)$ e $d\mid a,b$. Se $d=\alpha a+\beta b$ per certi interi $\alpha,\beta$, allora $d$ è un massimo comune divisore. Se $d>0$, è $\operatorname{MCD}(a,b)$.
\end{proposizione}
\begin{proof}
La condizione di divisore comune è già un'ipotesi. Se $c\mid a,b$, la proposizione~\ref{prop:combinazioni} dà $c\mid\alpha a+\beta b=d$. Vale quindi anche la condizione universale della definizione.
\end{proof}

\begin{definizione}
Gli interi $a,b$, non entrambi nulli, sono \textbf{coprimi} se $\operatorname{MCD}(a,b)=1$.
\end{definizione}

\begin{proposizione}
Se $d=\operatorname{MCD}(a,b)$, $a=da'$ e $b=db'$, allora $\operatorname{MCD}(a',b')=1$.
\end{proposizione}
\begin{proof}
Il numero $1$ divide $a',b'$. Per verificare che è un massimo comune divisore, sia $c$ un loro divisore comune. Poiché $a',b'$ non sono entrambi nulli, $c\ne0$. Il numero $dc$ divide $da'=a$ e $db'=b$, quindi $dc\mid d$. Esiste perciò $q\in\Z$ con $d=dcq$. Essendo $d>0$, si cancella e si ottiene $1=cq$, cioè $c=\pm1$. Ogni divisore comune divide dunque $1$, come richiesto.
\end{proof}

\section{Divisione euclidea}
% A–L L10 pp. 2–3; M–Z L10 pp. 1–3. Induzione comune resa esplicita in tutte le sue parti.
\begin{teorema}[Divisione euclidea]\label{teo:divisione}
Per $a,b\in\Z$, $b\ne0$, esistono unici $q,r\in\Z$ tali che
\[
a=bq+r,\qquad 0\le r<|b|.
\]
Si chiamano \textbf{quoziente} e \textbf{resto} della divisione di $a$ per $b$.
\end{teorema}
\begin{proof}
\textbf{Esistenza per $b>0$, $a\ge0$: base.} Fissiamo $b>0$ e procediamo per induzione su $a\in\N$. Per $a=0$ si scelgono $q=0$, $r=0$; il vincolo $0\le0<b$ vale.

\textbf{Ipotesi e passo induttivo.} Supponiamo che per un certo $a\ge0$ esistano $q,r$ con $a=bq+r$ e $0\le r<b$. Allora $a+1=bq+(r+1)$. Poiché $r$ è intero, $1\le r+1\le b$. Se $r+1<b$, la nuova coppia è $(q,r+1)$; se $r+1=b$, si riscrive $a+1=b(q+1)+0$, e la coppia è $(q+1,0)$. In entrambi i casi il resto è nell'intervallo richiesto. Il principio d'induzione dà l'esistenza per ogni $a\ge0$.

\textbf{Esistenza per $b>0$, $a<0$.} Si applica quanto provato a $-a=bq_1+r_1$, $0\le r_1<b$. Se $r_1=0$, $a=b(-q_1)+0$. Se $r_1>0$,
\[
a=-bq_1-r_1=b(-q_1-1)+(b-r_1),\qquad 0<b-r_1<b.
\]
Il quoziente va quindi diminuito di un'unità per mantenere il resto non negativo.

\textbf{Esistenza per $b<0$.} Dividiamo $a$ per $-b>0$: dalla scrittura $a=(-b)q_1+r$ si ottiene $a=b(-q_1)+r$, con $0\le r<|-b|=|b|$.

\textbf{Unicità.} Se $a=bq+r=bq'+r'$, allora $b(q-q')=r'-r$. I vincoli sui resti danno $|r'-r|<|b|$. Se $q\ne q'$, essendo interi si avrebbe $|q-q'|\ge1$, quindi $|b(q-q')|\ge|b|$, contraddizione. Dunque $q=q'$ e poi $r=r'$.
\end{proof}

\begin{esempio}[Segni del quoziente e resto]
Le divisioni
\[
\begin{aligned}
27&=12\cdot2+3,&-5&=12\cdot(-1)+7,\\
8&=(-3)\cdot(-2)+2,&-10&=(-3)\cdot4+2
\end{aligned}
\]
hanno tutte resto non negativo e minore del valore assoluto del divisore. Per $0\le a<b$ si ha direttamente $q=0$, $r=a$.
\end{esempio}

\begin{corollario}
Per $b\ne0$, $b\mid a$ se e solo se il resto è zero.
\end{corollario}
\begin{proof}
Un resto zero dà $a=bq$. Viceversa, se $a=bk$, la coppia $(k,0)$ soddisfa tutti i vincoli, quindi per unicità è la coppia della divisione euclidea.
\end{proof}

\begin{lemma}[Passo euclideo]\label{lem:passo-euclideo}
Se $a=bq+r$, i divisori comuni di $a,b$ sono esattamente quelli di $b,r$. Di conseguenza le due coppie hanno gli stessi massimi comuni divisori, quando definiti.
\end{lemma}
\begin{proof}
\passomath{c\mid a,b\Rightarrow c\mid b,r} Dall'identità $r=a-qb$ e dalla proprietà delle combinazioni segue $c\mid r$; $c\mid b$ è già noto.

\passomath{c\mid b,r\Rightarrow c\mid a,b} Da $a=bq+r$ segue $c\mid a$; $c\mid b$ è già noto. Coincidendo tutti i divisori comuni, coincidono anche le due condizioni che definiscono un massimo comune divisore.
\end{proof}

\section{Algoritmo di Euclide e Bézout}
% A–L L11 pp. 1–2; M–Z L11 pp. 2–4. Separati terminazione, invariante e ricorrenza dei coefficienti.
\begin{teorema}[Algoritmo di Euclide]\label{teo:euclide}
Per $(a,b)\ne(0,0)$ esiste il massimo comune divisore positivo. Sostituiti gli argomenti con $A=|a|$, $B=|b|$, se $B=0$ il risultato è $A$; altrimenti è l'ultimo termine non nullo nelle divisioni successive
\[
r_0=A,\quad r_1=B,\quad
r_{i-1}=q_i r_i+r_{i+1},\qquad 0\le r_{i+1}<r_i.
\]
\end{teorema}
\begin{proof}
\textbf{Terminazione.} Se $B=0$, il risultato è già noto. Sia $B>0$. A ogni passo con $r_i>0$ la divisione euclidea esiste e determina il resto. I termini $r_1,r_2,\ldots$, finché positivi, sono strettamente decrescenti. Se fossero tutti positivi, il loro insieme avrebbe per buon ordinamento un minimo, seguito da un termine ancora minore: impossibile. Esiste quindi $k\ge1$ con $r_k>0$ e $r_{k+1}=0$.

\textbf{Invariante, base.} Indichiamo con $\mathcal D_i$ l'insieme dei divisori comuni di $r_{i-1},r_i$, per $1\le i\le k+1$. Per $i=1$ sono esattamente i divisori comuni di $A,B$, quindi di $a,b$, perché il segno non modifica la divisibilità.

\textbf{Invariante, passo induttivo.} Supponiamo che $\mathcal D_i$ coincida con l'insieme dei divisori comuni di $a,b$, per un $i\le k$. Dalla divisione $r_{i-1}=q_i r_i+r_{i+1}$, il lemma~\ref{lem:passo-euclideo} dà $\mathcal D_{i+1}=\mathcal D_i$. La stessa proprietà vale dunque per la coppia successiva. Per induzione vale fino a $\mathcal D_{k+1}$.

\textbf{Identificazione finale.} La coppia finale è $(r_k,0)$, il cui massimo comune divisore positivo è $r_k$. Poiché ha esattamente gli stessi divisori comuni della coppia iniziale, $r_k$ soddisfa le due condizioni di massimo comune divisore anche per $a,b$. Questo dimostra insieme esistenza e correttezza dell'algoritmo. Se il primo resto è già zero, $k=1$ e il risultato è $r_1=B$.
\end{proof}

\subsection{Algoritmo esteso e identità di Bézout}
\begin{teorema}[Identità di Bézout]\label{teo:bezout}
Per $(a,b)\ne(0,0)$ esistono interi $\alpha,\beta$ tali che
\[
\operatorname{MCD}(a,b)=\alpha a+\beta b.
\]
I coefficienti si possono calcolare insieme alle divisioni euclidee.
\end{teorema}
\begin{proof}
\textbf{Inizializzazione.} Siano $A=|a|$, $B=|b|$. Se $B=0$, $A=\pm a$ dà subito la tesi. Per $B>0$, usiamo i resti dell'algoritmo precedente e definiamo
\[
s_0=1,\quad t_0=0,\qquad s_1=0,\quad t_1=1.
\]
Si hanno $r_0=s_0A+t_0B=A$ e $r_1=s_1A+t_1B=B$: sono le due righe iniziali della costruzione.

\textbf{Ricorrenza e passo induttivo.} Per ogni divisione $r_{i-1}=q_i r_i+r_{i+1}$ poniamo
\[
s_{i+1}=s_{i-1}-q_i s_i,\qquad t_{i+1}=t_{i-1}-q_i t_i.
\]
Supponiamo che $r_{i-1}=s_{i-1}A+t_{i-1}B$ e $r_i=s_iA+t_iB$. Sostituendo nella formula del resto,
\[
\begin{aligned}
r_{i+1}&=r_{i-1}-q_i r_i\\
&=(s_{i-1}-q_i s_i)A+(t_{i-1}-q_i t_i)B\\
&=s_{i+1}A+t_{i+1}B.
\end{aligned}
\]
La ricorrenza usa soltanto interi e conserva quindi coefficienti interi. Dalle due basi e dal passo segue per induzione che ogni riga verifica $r_i=s_iA+t_iB$.

\textbf{Ultima riga e segni.} All'ultimo termine non nullo, $d=\operatorname{MCD}(a,b)=r_k=s_k|a|+t_k|b|$. Scriviamo $|a|=\varepsilon_a a$, $|b|=\varepsilon_b b$, dove $\varepsilon_z=1$ per $z\ge0$ e $\varepsilon_z=-1$ per $z<0$. Allora $\alpha=s_k\varepsilon_a$ e $\beta=t_k\varepsilon_b$ sono i coefficienti richiesti rispetto agli argomenti originali.
\end{proof}

\begin{esempio}[Divisioni, coefficienti e sostituzione a ritroso]
Per $(a,b)=(72,420)$ le divisioni sono
\[
\begin{aligned}
72&=0\cdot420+72,\\
420&=5\cdot72+60,\\
72&=1\cdot60+12,\\
60&=5\cdot12+0.
\end{aligned}
\]
La tabella dell'algoritmo esteso conserva l'identità $r_i=s_i\cdot72+t_i\cdot420$. Il quoziente $q_i$ nella riga $i$ serve per ottenere il resto $r_{i+1}$:
\[
\begin{array}{c|r|r|r|r}
i&r_i&q_i&s_i&t_i\\\hline
0&72&-&1&0\\
1&420&0&0&1\\
2&72&5&1&0\\
3&60&1&-5&1\\
4&12&5&6&-1\\
5&0&-&-35&6
\end{array}
\]
L'ultima riga non nulla dà $d=12=6\cdot72-420$. Per esempio la riga $4$ nasce da $s_4=s_2-q_3s_3=1-1\cdot(-5)=6$ e $t_4=t_2-q_3t_3=0-1\cdot1=-1$.

Si possono ottenere gli stessi coefficienti risalendo nelle divisioni: $12=72-60=72-(420-5\cdot72)=6\cdot72-420$. Questo è il procedimento a ritroso, mentre la tabella costruisce i coefficienti in avanti. Per la coppia $(72,-420)$ il massimo comune divisore resta $12$, e l'identità diventa $12=6\cdot72+1\cdot(-420)$.
\end{esempio}

\begin{corollario}
Per $(a,b)\ne(0,0)$, gli interi $a,b$ sono coprimi se e solo se esistono $\alpha,\beta\in\Z$ con $\alpha a+\beta b=1$.
\end{corollario}
\begin{proof}
\passomath{\Rightarrow} Bézout, applicato a $\operatorname{MCD}(a,b)=1$, dà la combinazione richiesta.

\passomath{\Leftarrow} Il numero $1$ divide $a,b$ ed è una loro combinazione intera positiva. Il criterio~\ref{prop:criterio-mcd} dà $\operatorname{MCD}(a,b)=1$.
\end{proof}

\begin{esercizio}[Non unicità dei coefficienti]
Mostrare che i coefficienti di Bézout non sono unici.
\end{esercizio}
\begin{svolgimento}
Se $d=\alpha a+\beta b$, per ogni $t\in\Z$ vale
\[
d=(\alpha+tb)a+(\beta-ta)b,
\]
perché i termini aggiunti sono $tab-tba=0$. Poiché $(a,b)\ne(0,0)$, almeno uno dei due coefficienti cambia quando cambia $t$: si ottengono infinite coppie distinte. Questa costruzione prova la non unicità, senza pretendere di descrivere qui tutte le coppie possibili.
\end{svolgimento}

\begin{corollario}
Per $d=\operatorname{MCD}(a,b)$, $(a,b)\ne(0,0)$,
\[
\{\alpha a+\beta b\mid\alpha,\beta\in\Z\}=d\Z,
\qquad d\Z=\{dk\mid k\in\Z\}.
\]
\end{corollario}
\begin{proof}
\passomath{\subseteq} Il numero $d$ divide $a,b$, quindi divide ogni loro combinazione intera.

\passomath{\supseteq} Bézout dà $d=\alpha_0a+\beta_0b$. Ogni $dk$ è allora $(k\alpha_0)a+(k\beta_0)b$, quindi appartiene all'insieme delle combinazioni.
\end{proof}

\subsection{La minima combinazione positiva}
% A–L L11 pp. 2–3; M–Z L10 pp. 3–4. Scelta la prova M–Z, indipendente da Bézout.
\begin{proposizione}\label{prop:min-combinazione}
Per $(a,b)\ne(0,0)$,
\[
\operatorname{MCD}(a,b)=\min\{\alpha a+\beta b>0\mid\alpha,\beta\in\Z\}.
\]
\end{proposizione}
\begin{proof}
\textbf{Esistenza del minimo.} L'insieme $S$ delle combinazioni positive è non vuoto: se $a\ne0$ contiene $|a|=\pm a+0b$; altrimenti contiene $|b|=0a\pm b$. Per buon ordinamento ha un minimo $s=\alpha_0a+\beta_0b>0$.

\textbf{Divisibilità di $a$.} Dividiamo $a$ per $s$: $a=qs+r$, $0\le r<s$. Allora
\[
r=a-qs=(1-q\alpha_0)a-q\beta_0b.
\]
Se $r>0$, sarebbe una combinazione positiva minore del minimo $s$. Quindi $r=0$ e $s\mid a$.

\textbf{Divisibilità di $b$.} Dividiamo anche $b=q's+r'$, $0\le r'<s$. Il resto è $r'=-q'\alpha_0a+(1-q'\beta_0)b$. Se positivo contraddirebbe la stessa minimalità, dunque $r'=0$ e $s\mid b$.

\textbf{Condizione universale.} Ogni divisore comune $c$ di $a,b$ divide la combinazione $s=\alpha_0a+\beta_0b$. Quindi $s$ soddisfa entrambe le condizioni di massimo comune divisore ed è positivo: è $\operatorname{MCD}(a,b)$.
\end{proof}

La prova mediante il minimo fornisce anche un'esistenza del massimo comune divisore e dell'identità di Bézout indipendente dall'algoritmo. L'algoritmo esteso offre in più un procedimento di calcolo dei coefficienti.

\section{Lemma di Euclide e fattorizzazione}
\begin{lemma}[Euclide]\label{lem:euclide}
Se $\operatorname{MCD}(a,b)=1$ e $a\mid bc$, allora $a\mid c$.
\end{lemma}
\begin{proof}
Da Bézout si ha $\alpha a+\beta b=1$. Moltiplicando per $c$, $c=\alpha ac+\beta bc$. Se $bc=am$, questa identità diventa $c=a(\alpha c+\beta m)$: il quoziente è intero, quindi $a\mid c$.
\end{proof}

\begin{corollario}
Se $a\mid c$, $b\mid c$ e $\operatorname{MCD}(a,b)=1$, allora $ab\mid c$.
\end{corollario}
\begin{proof}
% M–Z L11 p. 5: prova diretta via Bézout, senza divisione per a o b.
Scriviamo $c=au=bv$ e $1=\alpha a+\beta b$. Moltiplicando per $c$ e sostituendo nelle due parti,
\[
c=\alpha ac+\beta bc=\alpha abv+\beta bau=ab(\alpha v+\beta u).
\]
Il quoziente esibito è intero. L'argomento vale anche nei casi nulli consentiti dalla coprimalità.
\end{proof}

\begin{proposizione}[Primi e divisibilità dei prodotti]\label{prop:primo-prodotto}
Per $p\in\Z\setminus\{0,1,-1\}$ sono equivalenti: $p$ è primo; per ogni $a,b\in\Z$, da $p\mid ab$ segue $p\mid a$ oppure $p\mid b$.
\end{proposizione}
\begin{proof}
\passomath{\Rightarrow} Sia $p$ primo e $p\mid ab$. Se $p\mid a$ abbiamo già concluso. Altrimenti $d=\operatorname{MCD}(p,a)$ è un divisore positivo di $p$, dunque è $1$ o $|p|$. Il secondo caso implicherebbe $p\mid a$, escluso. Quindi $d=1$ e il lemma di Euclide dà $p\mid b$.

\passomath{\Leftarrow} Supponiamo la proprietà e sia $p=uv$. Allora $p\mid uv$, perciò $p\mid u$ oppure $p\mid v$. Nel primo caso $u=pk$ dà $p=pkv$: cancellando $p\ne0$, $kv=1$, quindi $v=\pm1$. Nel secondo caso, allo stesso modo $u=\pm1$. È la definizione di primo.
\end{proof}

\begin{osservazione}
Le esclusioni $p\ne0,\pm1$ sono necessarie: le unità dividono ogni intero, mentre per $p=0$ la proprietà sui prodotti segue dall'assenza di divisori dello zero. Nessuno di questi tre numeri è primo.
\end{osservazione}

\begin{teorema}[Fondamentale dell'aritmetica]\label{teo:fattorizzazione}
Ogni intero $n\ne0,\pm1$ si scrive come
\[
n=\varepsilon p_1^{a_1}\cdots p_r^{a_r},\qquad
\varepsilon\in\{1,-1\},\quad a_i\in\N^*,
\]
con primi positivi distinti $p_1,\ldots,p_r$. Il segno è quello di $n$ e i primi con i rispettivi esponenti sono unici a meno dell'ordine.
\end{teorema}
\begin{proof}
\textbf{Esistenza: base.} Basta trattare $n\ge2$ ed estrarre poi il segno. Procediamo per induzione forte su $n$. Per $n=2$, la scrittura è il solo primo $2$.

\textbf{Esistenza: ipotesi e passo.} Sia $n>2$ e supponiamo che ogni intero $m$ con $2\le m<n$ sia prodotto di primi positivi. Se $n$ è primo, è già un prodotto con un fattore. Se non lo è, $n=uv$ con $1<u,v<n$. L'ipotesi si applica separatamente a $u$ e $v$; concatenando le due liste si ottiene una lista di primi di prodotto $n$. Il principio d'induzione forte conclude l'esistenza. Raccogliendo i fattori uguali si ottiene la forma con esponenti positivi.

\textbf{Unicità: parametro e base.} Consideriamo due scritture con ripetizioni ammesse,
\[
p_1\cdots p_\ell=q_1\cdots q_s,
\]
e proviamo per induzione su $\ell\ge0$ che $s=\ell$ e le liste coincidono a meno dell'ordine. Per $\ell=0$ il primo prodotto è $1$; un prodotto non vuoto di primi positivi è almeno $2$, quindi $s=0$.

\textbf{Unicità: passo.} Supponiamo la proprietà per $\ell-1$ e sia $\ell\ge1$. Il prodotto a sinistra è maggiore di $1$, dunque $s\ge1$. Poiché $p_1\mid q_1\cdots q_s$, applicazioni successive della proposizione~\ref{prop:primo-prodotto} mostrano che $p_1$ divide almeno un $q_j$. Infatti, se non divide il primo fattore, deve dividere il prodotto dei successivi, e si ripete fino all'ultimo. Essendo $q_j$ primo positivo e $p_1>1$, si ha $p_1=q_j$. Riordiniamo il secondo prodotto portando $q_j$ al primo posto e cancelliamo il fattore non nullo comune. Restano prodotti con $\ell-1$ e $s-1$ fattori. L'ipotesi induttiva dà $s-1=\ell-1$ e identifica tutte le occorrenze rimanenti.

\textbf{Esponenti e segno.} L'unicità delle liste con ripetizioni determina quante volte compare ciascun primo: sono gli esponenti. Per $n<0$ si applica il risultato a $|n|$ e si antepone il segno meno. Il segno è unico perché ogni prodotto di primi positivi è positivo.
\end{proof}

\subsection{Formule per MCD e mcm}
% A–L L11 p. 7: completati i passaggi mediante gli esponenti, senza aggiungere nuovi argomenti.
\begin{proposizione}\label{prop:mcd-mcm}
Per $a,b\ne0$ usiamo una stessa lista di primi positivi distinti, ammettendo esponenti nulli:
\[
a=\varepsilon\prod_{i=1}^r p_i^{\alpha_i},\qquad
b=\eta\prod_{i=1}^r p_i^{\beta_i},\qquad \alpha_i,\beta_i\in\N.
\]
Allora
\[
\operatorname{MCD}(a,b)=\prod_{i=1}^r p_i^{\min(\alpha_i,\beta_i)},\qquad
\operatorname{mcm}(a,b)=\prod_{i=1}^r p_i^{\max(\alpha_i,\beta_i)}.
\]
In particolare il minimo comune multiplo positivo esiste e $\operatorname{MCD}(a,b)\operatorname{mcm}(a,b)=|ab|$.
\end{proposizione}
\begin{proof}
\textbf{Divisibilità ed esponenti.} Per interi non nulli $u,v$, $u\mid v$ se e solo se ogni primo ha in $u$ esponente non maggiore che in $v$. Se $v=uw$, gli esponenti si sommano per l'unicità della fattorizzazione, quindi la disuguaglianza è necessaria. Viceversa le differenze non negative degli esponenti costruiscono un intero $w$ con $v=uw$, scegliendo il segno opportuno.

\textbf{Massimo comune divisore.} Il prodotto $d$ formato usando gli esponenti minimi divide $a$ e $b$. Ogni divisore comune $c$ è non nullo, ha solo primi presenti in entrambi e per ciascuno ha esponente al più il minimo dei due esponenti: quindi $c\mid d$. Le due condizioni della definizione sono soddisfatte e $d>0$.

\textbf{Minimo comune multiplo.} Il prodotto $m$ formato usando gli esponenti massimi è multiplo di $a$ e $b$. Ogni altro multiplo comune non nullo deve avere esponenti almeno pari a quei massimi, quindi è divisibile per $m$. Anche il multiplo comune $0$ è divisibile per $m$. Si ottiene la definizione e $m>0$.

\textbf{Prodotto.} Per ogni $i$, $\min(\alpha_i,\beta_i)+\max(\alpha_i,\beta_i)=\alpha_i+\beta_i$. Moltiplicando le due formule si ottiene quindi $|a||b|=|ab|$.
\end{proof}

I prodotti vuoti valgono $1$: le formule includono così $a=\pm1$ o $b=\pm1$. Se un solo argomento è nullo, l'identità del prodotto rimane valida con $\operatorname{MCD}(a,0)=|a|$ e $\operatorname{mcm}(a,0)=0$.

% Ripresa M–Z L12 p. 2: conseguenza operativa della formula già dimostrata.
Per $a,b\ne0$ si può dunque calcolare anche il minimo comune multiplo mediante l'algoritmo di Euclide, senza fattorizzare gli argomenti:
\[
\operatorname{mcm}(a,b)=\frac{|ab|}{\operatorname{MCD}(a,b)}.
\]

\begin{esercizio}[Radice di un primo]
Dimostrare che per ogni primo positivo $p$, $\sqrt p\notin\Q$.
\end{esercizio}
\begin{svolgimento}
Supponiamo $\sqrt p=a/b$ con $a\in\Z$, $b\in\N^*$. Dividendo entrambi per $\operatorname{MCD}(a,b)$ possiamo scegliere $a,b$ coprimi. Dall'uguaglianza $a^2=pb^2$ segue $p\mid a^2$, quindi $p\mid a$ per la proprietà sui prodotti. Scriviamo $a=pc$: sostituendo e cancellando $p\ne0$, si ottiene $pc^2=b^2$. Allora $p\mid b^2$, quindi $p\mid b$. Ma un primo $p>1$ non può dividere entrambi gli interi coprimi $a,b$. La contraddizione esclude la razionalità.

% M–Z L12 p. 3: metodo alternativo mediante gli esponenti, senza frazione ridotta.
\textbf{Secondo metodo.} Dalla stessa uguaglianza $a^2=pb^2$, con $a,b\ne0$, confrontiamo l'esponente del primo $p$ nelle due fattorizzazioni. Se $p$ compare in $|a|$ con esponente $h\ge0$ e in $|b|$ con esponente $k\ge0$, nell'uguaglianza dovrebbe valere $2h=1+2k$. Il primo membro è pari e il secondo dispari: impossibile. Questo argomento usa l'unicità della fattorizzazione e non richiede che $a,b$ siano coprimi.
\end{svolgimento}

% !TeX root = ../main.tex
\chapter{Aritmetica modulare}\label{cap:modulare}
\fonti{Lezione 12}

\section{Congruenza e classi resto}
% A–L L12 pp. 1–2; M–Z L12 pp. 4–5.
Fissiamo un intero $m\ge2$, detto \textbf{modulo}.
\begin{definizione}
Due interi $x,y$ sono \textbf{congrui modulo $m$} se hanno lo stesso resto nella divisione euclidea per $m$. Si scrive $x\equiv_m y$, oppure $x\equiv y\pmod m$.
\end{definizione}

\begin{proposizione}\label{prop:congruenza-divisibilita}
Per ogni $x,y\in\Z$,
\[
x\equiv y\pmod m\quad\Longleftrightarrow\quad m\mid x-y.
\]
\end{proposizione}
\begin{proof}
Scriviamo le divisioni $x=mq_1+r_1$, $y=mq_2+r_2$, con $0\le r_1,r_2<m$. Sottraendo,
\[
x-y=m(q_1-q_2)+(r_1-r_2).
\]
Se $r_1=r_2$, la differenza è multipla di $m$. Viceversa, se $m\mid x-y$, allora $m\mid r_1-r_2$. Poiché $|r_1-r_2|<m$, tale multiplo di $m$ deve essere zero: $r_1=r_2$.
\end{proof}

\begin{proposizione}
La congruenza modulo $m$ è una relazione di equivalenza su $\Z$.
\end{proposizione}
\begin{proof}
È riflessiva perché $m\mid x-x=0$. È simmetrica perché da $x-y=mk$ segue $y-x=m(-k)$. È transitiva perché da $x-y=mk$ e $y-z=mh$ segue $x-z=m(k+h)$.
\end{proof}

La stessa equivalenza è quella indotta dalla funzione resto
\[
\rho_m:\Z\longrightarrow\{0,1,\ldots,m-1\},\qquad
\rho_m(x)=\text{resto della divisione di $x$ per $m$}:
\]
due interi sono equivalenti esattamente quando hanno la stessa immagine.

\begin{definizione}
La \textbf{classe resto} di $x$ modulo $m$ è
\[
[x]_m=\{y\in\Z\mid m\mid x-y\}
     =\{x+mk\mid k\in\Z\}=x+m\Z.
\]
L'insieme quoziente si indica con $\Z/m\Z$. Quando il modulo è chiaro, scriviamo semplicemente $[x]$; si usa anche la notazione $\overline{x}$.
\end{definizione}

L'uguaglianza delle due descrizioni della classe segue da $m\mid x-y\iff y=x+mk$ per qualche $k\in\Z$.

\begin{proposizione}
Ogni classe ha un unico rappresentante in $\{0,1,\ldots,m-1\}$. In particolare
\[
\Z/m\Z=\{[0]_m,[1]_m,\ldots,[m-1]_m\},\qquad |\Z/m\Z|=m.
\]
\end{proposizione}
\begin{proof}
La divisione $x=mq+r$, $0\le r<m$, dà $[x]_m=[r]_m$, quindi ogni classe ha un tale rappresentante. Se $r,s$ sono entrambi in quell'intervallo e $[r]_m=[s]_m$, allora $m\mid r-s$ e $|r-s|<m$, dunque $r=s$. Le $m$ classi elencate sono perciò distinte ed esauriscono il quoziente.
\end{proof}

\begin{esempio}
Per $m=2$, $[0]_2=2\Z$ è l'insieme dei pari e $[1]_2=1+2\Z$ quello dei dispari. Per $m=3$ le classi sono $3\Z$, $1+3\Z$ e $2+3\Z$: per esempio $[-3]_3=[0]_3=[6]_3$ e $[8]_3=[5]_3=[2]_3$.
\end{esempio}

\begin{schema}
\begin{tikzpicture}[x=1.7cm,y=.65cm]
\draw[asse] (-2.35,0)--(2.5,0);
\foreach \k/\testo in {-2/{x-2m},-1/{x-m},0/{x},1/{x+m},2/{x+2m}}{
  \fill[coloreschema] (\k,0) circle (2pt);
  \node[below=5pt] at (\k,0) {$\testo$};
}
\foreach \k in {-2,-1,0,1}{\draw[coloreschema] (\k,.15) to[bend left=30] (\k+1,.15);}
\node at (0,1.1) {Rappresentanti della stessa classe $[x]_m$};
\end{tikzpicture}
\end{schema}

\begin{osservazione}
Tra classi si scrive un'\emph{uguaglianza}; tra interi, una \emph{congruenza}. Per esempio $[-2]_3=[1]_3$ e $-2\equiv1\pmod3$, mentre $-2\ne1$ in $\Z$.
\end{osservazione}

\section{Operazioni sul quoziente}
% A–L L12 pp. 3–4: scelta la prova esplicita, lasciata per esercizio in M–Z p. 5.
\begin{proposizione}\label{prop:operazioni-modulari}
Le regole
\[
[a]_m+[b]_m=[a+b]_m,\qquad [a]_m[b]_m=[ab]_m
\]
definiscono operazioni su $\Z/m\Z$: il risultato non dipende dai rappresentanti scelti.
\end{proposizione}
\begin{proof}
Siano $[a]_m=[a']_m$ e $[b]_m=[b']_m$, cioè $m\mid a-a'$ e $m\mid b-b'$. Per la somma,
\[
(a+b)-(a'+b')=(a-a')+(b-b')
\]
è divisibile per $m$, quindi $[a+b]_m=[a'+b']_m$. Per il prodotto,
\[
ab-a'b'=a(b-b')+(a-a')b'
\]
è ancora divisibile per $m$, quindi $[ab]_m=[a'b']_m$. Entrambe le operazioni sono dunque ben definite.
\end{proof}

\begin{proposizione}
Le operazioni sul quoziente sono associative e commutative; il prodotto è distributivo rispetto alla somma. I neutri sono $[0]_m$ e $[1]_m$, rispettivamente; ogni $[a]_m$ ha opposto $[-a]_m$. Pertanto $\Z/m\Z$ è un anello commutativo con unità.
\end{proposizione}
\begin{proof}
Omettiamo il pedice $m$. Le proprietà si verificano sui rappresentanti usando quelle degli interi:
\[
\begin{aligned}
[a]+[b]&=[a+b]=[b+a]=[b]+[a],\\
([a]+[b])+[c]&=[(a+b)+c]=[a+(b+c)]=[a]+([b]+[c]),\\
[a][b]&=[ab]=[ba]=[b][a],\\
([a][b])[c]&=[(ab)c]=[a(bc)]=[a]([b][c]),\\
[a]([b]+[c])&=[a(b+c)]=[ab+ac]=[a][b]+[a][c].
\end{aligned}
\]
Inoltre $[a]+[0]=[a]$, $[a]+[-a]=[0]$ e $[a][1]=[a]$. La distributività dall'altro lato segue dalla commutatività del prodotto. Infine $[0]\ne[1]$, perché $m\ge2$ non divide $1$.
\end{proof}

\begin{esempio}[Parità e orologio]
In $\Z/2\Z$, $[0]+[0]=[0]$, $[0]+[1]=[1]$ e $[1]+[1]=[0]$: sono le regole della parità. Inoltre $[1]=[-1]=[-3]$, perciò una classe può coincidere con il proprio opposto.

In $\Z/12\Z$ si calcola come su un orologio con lo zero al posto del dodici: ad esempio $[9]+[5]=[14]=[2]$. L'opposto di $[7]_{12}$ è $[-7]_{12}=[5]_{12}$.
\end{esempio}

\begin{schema}
\begin{tikzpicture}[scale=.85]
\draw[spazio] (0,0) circle (1.45);
\foreach \k in {0,...,11}{
  \fill[coloreschema] ({90-30*\k}:1.45) circle (1.5pt);
  \node at ({90-30*\k}:1.8) {$\k$};
}
\draw[-{Stealth},coloreschema,line width=.9pt] (90:1.15) arc (90:-60:1.15);
\node at (0,-2.35) {Le dodici classi di $\Z/12\Z$};
\end{tikzpicture}
\end{schema}

\section{Invertibili e divisori dello zero}
% A–L L12 pp. 5–6; integrazioni M–Z L12 p. 6.
\begin{esempio}
In $\Z/4\Z$ si ha $[2]+[2]=[0]$ e $[2][2]=[0]$, pur essendo $[2]\ne[0]$. In $\Z/6\Z$, $[2][3]=[0]$ con entrambi i fattori non nulli. L'assenza di divisori dello zero in $\Z$ non passa quindi a tutti i quozienti.

In $\Z/5\Z$, invece, $[2][3]=[6]=[1]$: le classi $[2]$ e $[3]$ sono inverse l'una dell'altra. Le altre classi non nulle sono $[1]$, inversa di sé, e $[4]=[-1]$, anch'essa inversa di sé. Anche in $\Z/2\Z$ ogni classe non nulla è invertibile.
\end{esempio}

\begin{definizione}
Una classe $[a]_m$ è \textbf{invertibile}, o è un'\textbf{unità}, se esiste $[b]_m$ tale che $[a]_m[b]_m=[1]_m$. Tale classe $[b]_m$ si chiama \textbf{inverso} di $[a]_m$.

Seguendo la convenzione della lezione, $[a]_m$ è un \textbf{divisore dello zero} se esiste una classe $[b]_m\ne[0]_m$ con $[a]_m[b]_m=[0]_m$. Questa definizione include $[0]_m$.
\end{definizione}

\begin{osservazione}
La classe nulla rientra nella definizione perché $[0]_m[1]_m=[0]_m$ e $[1]_m\ne[0]_m$. Quando si richiede anche $[a]_m\ne[0]_m$, si parla di divisore dello zero \emph{non nullo}; questa condizione aggiuntiva va tenuta presente confrontando convenzioni diverse.
\end{osservazione}

L'inverso, quando esiste, è unico: se $[a][b]=[1]=[a][c]$, allora
\[
[b]=[b][1]=[b]([a][c])=([b][a])[c]=[1][c]=[c].
\]
Lo indichiamo con $[a]_m^{-1}$. Per esempio $[5]_{12}^{-1}=[5]_{12}$, perché $25\equiv1\pmod{12}$.

\begin{esempio}
In $\Z/6\Z$, $[1]$ e $[5]=[-1]$ sono invertibili. Le classi $[0],[2],[3],[4]$ sono divisori dello zero: oltre al prodotto $[2][3]=[0]$, si ha $[4][3]=[12]=[0]$. Qui ogni classe è un'unità oppure un divisore dello zero. Il prossimo risultato mostra che ciò vale per ogni modulo.
\end{esempio}

\Needspace{17\baselineskip}
\subsection{Il criterio del massimo comune divisore}
% A–L L12 pp. 6–8; M–Z L12 p. 6 per la costruzione dell'inverso con Bézout.
\begin{teorema}\label{teo:unita-modulari}
Siano $m\ge2$ e $a\in\Z$. Allora:
\begin{enumerate}
\item $[a]_m$ è invertibile se e solo se $\operatorname{MCD}(a,m)=1$;
\item $[a]_m$ è divisore dello zero, nella convenzione che include $[0]_m$, se e solo se $\operatorname{MCD}(a,m)>1$.
\end{enumerate}
In particolare ogni classe appartiene esattamente a una delle due categorie.
\end{teorema}
\begin{proof}
\textbf{Invertibilità: necessità.} Se $[a]_m[b]_m=[1]_m$, esiste $q\in\Z$ con $ab-1=mq$. Quindi $1=ab-mq$ è una combinazione intera di $a,m$. Ogni divisore comune divide $1$, perciò $\operatorname{MCD}(a,m)=1$.

\textbf{Invertibilità: sufficienza e costruzione.} Se $\operatorname{MCD}(a,m)=1$, Bézout dà $\alpha a+\beta m=1$. Passando alle classi,
\[
[\alpha]_m[a]_m+[\beta]_m[m]_m=[1]_m.
\]
Poiché $[m]_m=[0]_m$, segue $[\alpha]_m[a]_m=[1]_m$: dunque $[a]_m^{-1}=[\alpha]_m$.

\textbf{Divisori dello zero: esclusione del caso coprimo.} Se $\operatorname{MCD}(a,m)=1$ e $[a]_m[b]_m=[0]_m$, allora $m\mid ab$. Il lemma di Euclide applicato a $m,a$ dà $m\mid b$, cioè $[b]_m=[0]_m$. Non esiste quindi un annullatore non nullo.

\textbf{Divisori dello zero: costruzione dell'annullatore.} Se $d=\operatorname{MCD}(a,m)>1$, scriviamo $a=da'$ e $m=dm'$. Si ha $m'=m/d\in\Z$ e $1\le m'<m$, quindi $[m']_m\ne[0]_m$. Inoltre
\[
[a]_m[m']_m=[da'm']_m=[ma']_m=[0]_m.
\]
Quindi $[a]_m$ è divisore dello zero, con annullatore non nullo $[m/d]_m$. La costruzione include $a\equiv0\pmod m$: in quel caso $d=m$ e l'annullatore è $[1]_m$.
\end{proof}

\begin{osservazione}
Il criterio non dipende dal rappresentante: se $a'=a+km$, un intero divide $a,m$ se e solo se divide $a',m$, mediante le combinazioni $a'=a+km$ e $a=a'-km$. I due massimi comuni divisori positivi coincidono.

Se $m=p$ è primo, ogni classe non nulla ha $\operatorname{MCD}(a,p)=1$ ed è invertibile. Dunque $\Z/p\Z$ è un \textbf{campo}: un anello commutativo con $[0]\ne[1]$ in cui ogni elemento non nullo è invertibile. Questo spiega gli esempi dei moduli $2$ e $5$.
\end{osservazione}

\subsection{Cancellazione e calcolo degli inversi}
\begin{osservazione}
Si può cancellare un fattore invertibile: moltiplicando $[a][b]=[a][c]$ per $[a]^{-1}$ si ottiene $[b]=[c]$. Un divisore dello zero, invece, può impedire la cancellazione. In $\Z/6\Z$,
\[
[2][1]=[2]=[2][4],\qquad [1]\ne[4].
\]
Dalla sola uguaglianza del prodotto a zero non segue che uno dei fattori sia nullo: $[2][3]=[0]$ è già un controesempio.
\end{osservazione}

\Needspace{10\baselineskip}
\begin{esercizio}[Calcoli modulo 84]
In $\Z/84\Z$, caratterizzare le classi invertibili e i divisori dello zero. Trovare una classe non nulla che annulli $[30]_{84}$ e calcolare $[65]_{84}^{-1}$.
\end{esercizio}
\begin{svolgimento}
Poiché $84=2^2\cdot3\cdot7$, una classe $[a]_{84}$ è divisore dello zero se e solo se almeno uno tra $2,3,7$ divide $a$. È invertibile se nessuno dei tre divide $a$.

Si ha $\operatorname{MCD}(30,84)=6$. La costruzione del teorema dà l'annullatore $[84/6]_{84}=[14]_{84}$, non nullo, e infatti $30\cdot14=420=5\cdot84$. Anche $[28]_{84}$ e $[42]_{84}$ sono annullatori non nulli, perché $30\cdot28=10\cdot84$ e $30\cdot42=15\cdot84$. Ogni loro multiplo annulla ancora $[30]$, ma il multiplo scelto può essere la classe nulla.

Per $65$ e $84$ le divisioni sono
\[
\begin{aligned}
84&=65+19,&65&=3\cdot19+8,\\
19&=2\cdot8+3,&8&=2\cdot3+2,\\
3&=2+1,&2&=2\cdot1.
\end{aligned}
\]
L'ultimo resto non nullo è $1$. Risalendo,
\[
\begin{aligned}
1&=3-2=3-(8-2\cdot3)=3\cdot3-8\\
 &=3(19-2\cdot8)-8=3\cdot19-7\cdot8\\
 &=3\cdot19-7(65-3\cdot19)=24\cdot19-7\cdot65\\
 &=24(84-65)-7\cdot65=24\cdot84-31\cdot65.
\end{aligned}
\]
Quindi $[65]_{84}^{-1}=[-31]_{84}=[53]_{84}$. La verifica diretta è $65\cdot53=3445=41\cdot84+1$.
\end{svolgimento}

\Needspace{9\baselineskip}
\begin{esercizio}
Trovare l'inverso di $[5]_{84}$. Stabilire se $[8]_{84}$ è invertibile.
\end{esercizio}
\begin{svolgimento}
Le divisioni $84=16\cdot5+4$ e $5=4+1$ danno
\[
1=5-4=5-(84-16\cdot5)=17\cdot5-84.
\]
Pertanto $[5]_{84}^{-1}=[17]_{84}$. Invece $\operatorname{MCD}(8,84)=4>1$, quindi $[8]_{84}$ non è invertibile. È un divisore dello zero: ad esempio $[8]_{84}[21]_{84}=[168]_{84}=[0]_{84}$ e $[21]_{84}\ne[0]_{84}$.
\end{svolgimento}

\section{Calcoli con le classi resto}
% A–L L12 p. 9.
\begin{esempio}[Il numero 2026]
In $\Z/10\Z$, $2026=202\cdot10+6$ dà $[2026]_{10}=[6]_{10}$.

In $\Z/3\Z$, poiché $[10]_3=[1]_3$, per le proprietà del prodotto si ha anche $[1000]_3=[10]_3^3=[1]_3$. Quindi
\[
\begin{aligned}
[2026]_3&=[2\cdot1000+2\cdot10+6]_3\\
&=[2]_3[1]_3+[2]_3[1]_3+[0]_3\\
&=[4]_3=[1]_3.
\end{aligned}
\]
Equivalentemente, $[2]_3=[-1]_3$ permette di ottenere $[-1]_3+[-1]_3=[-2]_3=[1]_3$. Il resto della divisione di $2026$ per $3$ è perciò $1$.
\end{esempio}

\end{document}
